\documentclass[12pt,a4paper,UTF8]{ctexart}
\usepackage[a4paper,top=22mm,bottom=22mm,left=22mm,right=22mm,headheight=15pt]{geometry}
\usepackage{amsmath,amssymb,mathtools,bm}
\usepackage{graphicx,xcolor,booktabs,longtable,array,enumitem,fancyhdr,caption,float,microtype,setspace}
\newcommand{\Needspace}[1]{\par\begingroup\dimen0=#1\relax\dimen2=\pagegoal\advance\dimen2 by-\pagetotal\ifdim\dimen0>\dimen2\newpage\fi\endgroup}
\newcommand{\needspace}[1]{\par\begingroup\skip0=\lastskip\vskip-\skip0\dimen0=#1\relax\advance\dimen0 by\skip0\vskip\dimen0\penalty-100\vskip-\dimen0\vskip\skip0\penalty9999\endgroup}
\usepackage[hidelinks,unicode]{hyperref}
\usepackage[most]{tcolorbox}
\definecolor{ContestBlue}{HTML}{17365D}
\hypersetup{pdftitle={第 047 套},pdfauthor={}}
\setstretch{1.24}
\setlength{\parindent}{2em}
\setlength{\parskip}{0.35em}
\setlength{\emergencystretch}{2em}
\setlist{itemsep=0.35em,topsep=0.4em}
\captionsetup{font=small,labelfont=bf,justification=centering}
\setcounter{secnumdepth}{0}
\pagestyle{fancy}
\fancyhf{}
\fancyhead[L]{\small 第 047 套}
\fancyfoot[C]{\thepage}
\renewcommand{\headrulewidth}{0.4pt}
\ctexset{
  section={format=\Large\bfseries\color{ContestBlue},beforeskip=1.2em,afterskip=0.7em},
  subsection={format=\large\bfseries\color{ContestBlue},beforeskip=1.0em,afterskip=0.5em},
  subsubsection={format=\normalsize\bfseries,beforeskip=0.8em,afterskip=0.35em}
}
\makeatletter
\newcommand{\tightlist}{\setlength{\itemsep}{0pt}\setlength{\parskip}{0pt}}
\newcommand{\maxwidth}{\ifdim\Gin@nat@width>\linewidth\linewidth\else\Gin@nat@width\fi}
\newcommand{\maxheight}{\ifdim\Gin@nat@height>0.72\textheight 0.72\textheight\else\Gin@nat@height\fi}
\newcommand{\pandocbounded}[1]{\begingroup\setkeys{Gin}{width=\maxwidth,height=\maxheight,keepaspectratio}#1\endgroup}
\makeatother
\begin{document}
\begin{titlepage}
  \centering
  \vspace*{58mm}
  {\Huge\bfseries\color{ContestBlue} 第 047 套\par}
  \vfill
  \vspace*{28mm}
\end{titlepage}
\section{第 1 题}\label{ux7b2c-1-ux9898}

四根完全相同的均匀细杆首尾以光滑铰链连接，每根长度
\(l=0.5\,\mathrm m\)、质量 \(m=0.1\,\mathrm{kg}\)。双手分别抓住两端
\(A,E\)，保持两手等高且水平距离为
\(L=1.6\,\mathrm m\)。系统静止时，最低铰点 \(C\) 与双手的高度差为多少？

\begin{figure}[H]
\centering
\includegraphics[width=0.8\linewidth,height=\textheight,keepaspectratio,alt={图1-1　四节折叠尺的对称平衡构型}]{../figures/question-01.pdf}
\caption*{图1-1　四节折叠尺的对称平衡构型}
\end{figure}

\subsection{一、题型定位}\label{ux4e00ux9898ux578bux5b9aux4f4d}

本题是多刚体铰接系统的势能极小问题。铰链光滑，静止构型应在几何约束下使总重力势能取极小值。左右完全对称，可用外侧杆、内侧杆的两个倾角描述。

\subsection{二、知识点}\label{ux4e8cux77e5ux8bc6ux70b9}

\begin{enumerate}
\def\labelenumi{\arabic{enumi}.}
\tightlist
\item
  对称性降维；
\item
  均匀杆质心位置；
\item
  约束极值与隐函数求导；
\item
  稳定平衡的势能判据；
\item
  三角函数数值求根。
\end{enumerate}

\subsection{三、物理图像}\label{ux4e09ux7269ux7406ux56feux50cf}

设外侧杆 \(AB,DE\) 相对水平线向下的角为 \(\theta_1\)，内侧杆 \(BC,CD\)
的角为
\(\theta_2\)。左右两半互为镜像。端点水平距离限制两角余弦之和，而总势能由四根杆质心的下降量决定。

\subsection{四、解题入口}\label{ux56dbux89e3ux9898ux5165ux53e3}

水平投影约束为 \[
2l(\cos\theta_1+\cos\theta_2)=L. \tag{1}
\] 最低点下降量为 \[
h=l(\sin\theta_1+\sin\theta_2). \tag{2}
\] 必须先由势能极小确定两角关系，不能仅凭图形假设四根杆倾角相等。

\subsection{五、解析}\label{ux4e94ux89e3ux6790}

取双手高度为零势能面。两根外侧杆的质心各下降 \[
\frac l2\sin\theta_1.
\] 两根内侧杆的质心各下降 \[
l\sin\theta_1+\frac l2\sin\theta_2.
\] 所以总势能 \[
\begin{aligned}
U
&=-2mg\frac l2\sin\theta_1
-2mg\left(l\sin\theta_1+\frac l2\sin\theta_2\right)\\
&=-mgl(3\sin\theta_1+\sin\theta_2). 
\end{aligned} \tag{3}
\] 稳定平衡要求在式（1）约束下使 \[
3\sin\theta_1+\sin\theta_2
\] 最大。

对式（1）求微分： \[
-\sin\theta_1\,\mathrm d\theta_1
-\sin\theta_2\,\mathrm d\theta_2=0,
\] 故 \[
\frac{\mathrm d\theta_2}{\mathrm d\theta_1}
=-\frac{\sin\theta_1}{\sin\theta_2}. \tag{4}
\] 对式（3）中的无量纲部分求导并令零： \[
3\cos\theta_1
+\cos\theta_2\frac{\mathrm d\theta_2}{\mathrm d\theta_1}=0.
\] 代入式（4）： \[
3\cos\theta_1
-\cos\theta_2\frac{\sin\theta_1}{\sin\theta_2}=0.
\] 整理得 \[
\boxed{\tan\theta_1=3\tan\theta_2}. \tag{5}
\]

\begin{figure}[H]
\centering
\includegraphics[width=0.8\linewidth,height=\textheight,keepaspectratio,alt={图1-2　几何约束与势能极值的角度关系}]{../figures/solution-01.pdf}
\caption*{图1-2　几何约束与势能极值的角度关系}
\end{figure}

令 \[
x=\tan\theta_2,\qquad \tan\theta_1=3x.
\] 则 \[
\cos\theta_2=\frac1{\sqrt{1+x^2}},
\qquad
\cos\theta_1=\frac1{\sqrt{1+9x^2}}.
\] 由式（1） \[
\frac1{\sqrt{1+9x^2}}
+\frac1{\sqrt{1+x^2}}
=\frac{L}{2l}
=1.6. \tag{6}
\] 数值求正根： \[
x\simeq0.375534.
\] 于是 \[
\sin\theta_2=\frac{x}{\sqrt{1+x^2}},
\qquad
\sin\theta_1=\frac{3x}{\sqrt{1+9x^2}}.
\] 代入式（2）： \[
\begin{aligned}
h
&=0.5\left[
\frac{3x}{\sqrt{1+9x^2}}
+\frac{x}{\sqrt{1+x^2}}
\right]\\
&\simeq0.549720\,\mathrm m.
\end{aligned}
\] 故 \[
\boxed{h\simeq0.550\,\mathrm m}. \tag{7}
\]

\textbf{数值根的几何复核。} 由 \(x=0.375534\) 得 \[
\theta_2=\arctan x\simeq20.58^\circ,
\qquad
\theta_1=\arctan(3x)\simeq48.41^\circ.
\] 水平投影 \[
2l(\cos\theta_1+\cos\theta_2)
\simeq1.600\,{\rm m}=L,
\] 竖直投影为
\(l(\sin\theta_1+\sin\theta_2)\simeq0.550\,{\rm m}\)。两式同时成立才说明正根对应真实构型；只代回一个方程不足以排除错误支路。

\textbf{从三根内力建立平衡。} 四段均匀杆的质量同为
\(m\)。从最中间节点两侧的对称性可知，该节点的水平内力互相抵消。对左侧两杆分别取各自外端点的力矩，得到外段与内段所需的水平分力关系。也可把每根杆的重心下降量写为
\[
y_{\rm out}=\frac l2\sin\theta_1,\qquad
y_{\rm in}=l\sin\theta_1+\frac l2\sin\theta_2.
\] 两侧总势能取上端水平线为零： \[
U=-mgl(3\sin\theta_1+\sin\theta_2).
\] 长度约束 \[
\cos\theta_1+\cos\theta_2=\frac L{2l}=c_0.
\] 在固定 \(L\) 下
\(d\theta_2/d\theta_1=-\sin\theta_1/\sin\theta_2\)。于是 \[
\frac{dU}{d\theta_1}
=-mgl\left(3\cos\theta_1
-\cos\theta_2\frac{\sin\theta_1}{\sin\theta_2}\right).
\] 令零得
\(\tan\theta_1=3\tan\theta_2\)。系数三来自外杆重心下降与内杆整体下降的共同贡献，不是四杆数目直接给出。

\textbf{平衡根的唯一性。} 令 \(x=\tan\theta_2>0\)： \[
F(x)=\frac1{\sqrt{1+9x^2}}+\frac1{\sqrt{1+x^2}}.
\] 其导数 \[
F'(x)=-\frac{9x}{(1+9x^2)^{3/2}}
-\frac{x}{(1+x^2)^{3/2}}<0.
\] 且 \(F(0)=2,F(\infty)=0\)，所以对 \(0<L<4l\)
只有一个下垂根。不能为了凑角度用两个不满足长度约束的独立反三角值。

\textbf{跨度的敏感性。} 由 \(F(x)=L/(2l)\)： \[
\frac{dx}{dL}=\frac1{2lF'(x)}<0.
\] 下垂高度 \[
h(x)=l\left[\frac{3x}{\sqrt{1+9x^2}}
+\frac{x}{\sqrt{1+x^2}}\right],
\] \[
h'(x)=l\left[\frac3{(1+9x^2)^{3/2}}
+\frac1{(1+x^2)^{3/2}}\right]>0.
\] 因此
\(dh/dL<0\)。跨度增大，内外角和下垂高度都减小，与图1-2的几何一致。

\textbf{两端极限的展开。} 接近平直，令 \(\epsilon=2-L/(2l)\)： \[
F(x)=2-5x^2+\frac{123}4x^4+O(x^6).
\] 首阶 \(x\simeq\sqrt{\epsilon/5}\)，高度
\(h\simeq4l\sqrt{\epsilon/5}\)。下垂不是与跨度短缺成正比，而是平方根关系；因此接近平直时几何对长度误差比较敏感。

跨度趋零时 \(x\gg1\)： \[
F(x)\simeq\frac4{3x},\qquad
x\simeq\frac{8l}{3L},\qquad h\to2l.
\]
角度趋近直角，四杆折成两侧竖直链，支撑点同处一条竖线的极限存在，但真实铰链与杆厚会限制折叠范围。

\textbf{求双手的支持力。} 两侧对称，每只手的竖直支持力为
\(2mg\)。水平分力不能由整体受力直接求出，因为它们互相抵消。让端点距离作微小准静态变化，定义每只手向外的水平力为
\(H\)，两手的总功为 \(H\,dL\)。 \[
dU_{\rm eq}=H\,dL.
\] 在原来的长度约束中引入乘子 \(\Lambda\)： \[
\mathcal F=U+\Lambda[2l(\cos\theta_1+\cos\theta_2)-L].
\] 由两个角的一阶条件， \[
-3mgl\cos\theta_1-2\Lambda l\sin\theta_1=0,
\] \[
-mgl\cos\theta_2-2\Lambda l\sin\theta_2=0.
\] 所以 \[
\Lambda=-\frac{3mg}{2}\cot\theta_1
=-\frac{mg}{2}\cot\theta_2,\qquad
\boxed{H=-\Lambda=\frac{mg}{2x}}.
\] 取 \(g=9.8\,\mathrm{m\,s^{-2}}\)，得到
\(H\simeq1.305\,\mathrm N\)，单手合力大小为 \[
F_{\rm hand}=\sqrt{(2mg)^2+H^2}\simeq2.355\,\mathrm N.
\] 向外拉直时 \(H>0\)，与平衡势能随跨度增大的趋势一致。高度本身与
\(m,g\) 无关，支持力却与它们成正比。

\textbf{固定端点后的对称小振动。}
这是增加的变式：端点固定，允许各杆在竖直平面内运动，仅考察左右镜面对称的扰动。令
\(q=\theta_1\)，长度约束使 \(\theta_2=\theta_2(q)\)，在平衡点记 \[
b=\frac{d\theta_2}{dq}=-\frac{\sin\theta_1}{\sin\theta_2}.
\] 四杆的平动与转动动能相加为 \[
T=ml^2\left[
\frac43\dot\theta_1^2+\frac13\dot\theta_2^2
+\cos(\theta_1-\theta_2)\dot\theta_1\dot\theta_2
\right].
\]
其中第二杆质心的速度平方包含交叉项；把四根杆分别绕固定端点转动会遗漏它。于是约束后的有效惯量为
\[
\mathcal M(q)=2ml^2\left[
\frac43+\frac{b^2}3+b\cos(\theta_1-\theta_2)
\right].
\] 约束的二阶导数为 \[
b'=-\frac{\cos\theta_1}{\sin\theta_2}
-\frac{\sin^2\theta_1\cos\theta_2}{\sin^3\theta_2}.
\] 势能曲率因此是 \[
K=U''(q_*)=mgl\left[
3\sin\theta_1+\sin\theta_2\,b^2-\cos\theta_2\,b'
\right]_{q_*}>0.
\] 令 \(\xi=q-q_*\)，二阶拉格朗日量为 \[
\mathcal L^{(2)}=\frac12\mathcal M_*\dot\xi^2-\frac12K\xi^2,
\] 故 \[
\boxed{\omega_{\rm sym}^2=K/\mathcal M_*}.
\]
这个结果只描述所选的对称模。若释放反对称扰动，应再引入一个独立坐标并检查完整矩阵，不能把单模频率当成全部频谱。

\needspace{14\baselineskip}

\textbf{误差传播的数值例子。} 高度和跨度都由同一个根 \(x\) 给出。令
\(F(x)=L/(2l)\)，则 \[
\frac{dh}{dL}=
\frac{\displaystyle
3(1+9x^2)^{-3/2}+(1+x^2)^{-3/2}}
{2F'(x)}.
\] 分母为负，跨度偏大使高度偏小。若杆长固定而跨度测量误差为
\(\delta L\)，一阶有 \[
\delta h\simeq\frac{dh}{dL}\delta L.
\] 若改为所有长度一起按比例放大，则 \(L/l\) 不变，角度不变，高度按 \(l\)
成比例增加： \[
h=l\,\mathcal H(L/l),\qquad
dh=\left(\mathcal H-\frac Ll\mathcal H'\right)dl+\mathcal H'\,dL.
\]
因此``杆长误差''和``端点间距误差''不是同一种扰动。进行数值复核时应先说明哪些量保持不变，否则同一个公式的差分检查可能得到不同导数。

\newpage

\section{第 2 题}\label{ux7b2c-2-ux9898}

阁楼顶棚为直线斜面，最高点正下方地面为发射点；最高点高度
\(h=15\,\mathrm m\)，顶棚与水平面夹角 \(\alpha=30^\circ\)。从地面以
\(v_0=40\,\mathrm{m/s}\) 发射小物块，取 \(g=10\,\mathrm{m/s^2}\)。

（1）出射角 \(\theta=35^\circ\) 时，求首次击中顶棚的位置与速度；

（2）若 \(\theta\) 可调，求飞至顶棚时间最短的出射角。

\begin{figure}[H]
\centering
\includegraphics[width=0.792\linewidth,height=\textheight,keepaspectratio,alt={图2-1　抛体轨迹与倾斜顶棚}]{../figures/question-02.pdf}
\caption*{图2-1　抛体轨迹与倾斜顶棚}
\end{figure}

\subsection{一、题型定位}\label{ux4e00ux9898ux578bux5b9aux4f4d-1}

本题是抛体运动与直线边界的首次相交问题。第二问不是最大射程，而是在给定初速度下最大化沿顶棚法向的初速度分量，从而缩短首次相交时间。

\subsection{二、知识点}\label{ux4e8cux77e5ux8bc6ux70b9-1}

\begin{enumerate}
\def\labelenumi{\arabic{enumi}.}
\tightlist
\item
  平抛坐标方程；
\item
  斜面直线方程；
\item
  二次方程的较小正根；
\item
  速度分解；
\item
  三角函数振幅合成。
\end{enumerate}

\subsection{三、物理图像}\label{ux4e09ux7269ux7406ux56feux50cf-1}

取发射点为原点，\(x\) 水平向阁楼低端，\(y\) 竖直向上。顶棚方程为
\(y=h-x\tan\alpha\)。抛物线与顶棚一般有两个代数交点，实际首次撞击必须取较小正时间。

\subsection{四、解题入口}\label{ux56dbux89e3ux9898ux5165ux53e3-1}

抛体方程 \[
x=v_0\cos\theta\,t,\qquad
y=v_0\sin\theta\,t-\frac12gt^2. \tag{1}
\] 代入顶棚方程即可得到关于 \(t\) 的二次方程。

\subsection{五、解析}\label{ux4e94ux89e3ux6790-1}

相交条件 \[
v_0\sin\theta\,t-\frac12gt^2
=h-v_0\cos\theta\,t\tan\alpha.
\] 整理： \[
\frac12gt^2
-v_0(\sin\theta+\cos\theta\tan\alpha)t+h=0. \tag{2}
\] 令 \[
A=v_0(\sin35^\circ+\cos35^\circ\tan30^\circ)
\simeq41.85\,\mathrm{m/s}.
\] 式（2）的较小正根 \[
t_1=\frac{A-\sqrt{A^2-2gh}}g
\simeq0.375\,\mathrm s. \tag{3}
\]

\begin{figure}[H]
\centering
\includegraphics[width=0.664\linewidth,height=\textheight,keepaspectratio,alt={图2-2　首次交点的位置与速度分解}]{../figures/solution-02.pdf}
\caption*{图2-2　首次交点的位置与速度分解}
\end{figure}

坐标为 \[
x_1=40\cos35^\circ\,t_1
\simeq12.3\,\mathrm m,
\] \[
y_1=h-x_1\tan30^\circ
\simeq7.90\,\mathrm m.
\] 速度分量 \[
v_x=v_0\cos35^\circ\simeq32.8\,\mathrm{m/s},
\] \[
v_y=v_0\sin35^\circ-gt_1
\simeq19.2\,\mathrm{m/s}.
\] 故 \[
\boxed{\boldsymbol r_1\simeq(12.3,7.90)\,\mathrm m},
\] \[
\boxed{\boldsymbol v_1\simeq(32.8,19.2)\,\mathrm{m/s}}. \tag{4}
\]

第二问中式（2）在 \(h,g,v_0\) 固定时，首次相交时间随 \[
F(\theta)=\sin\theta+\tan\alpha\cos\theta
\] 增大而减小。利用 \[
F(\theta)=\frac{\sin(\theta+\alpha)}{\cos\alpha},
\] 最大值在 \[
\theta+\alpha=90^\circ
\] 取得。因此 \[
\boxed{\theta_{\min t}=90^\circ-\alpha=60^\circ}. \tag{5}
\]

\textbf{最短时间不是最高点条件。} 首次交顶棚时间中的角度因子为 \[
F(\theta)=\sin\theta+\tan\alpha\cos\theta
=\frac{\sin(\theta+\alpha)}{\cos\alpha}.
\] 最大化 \(F\) 才使时间最短，所以
\(\theta+\alpha=90^\circ\)。这不是抛体轨迹最高点条件；后者要求相交瞬间竖直速度为零，两者一般不同。

候选角还要检查交点位于实际屋面线段上，且所取是 \(t>0\)
的首次根。若屋面有限，内部极值落在屋面之外时，应改在端点比较。

\textbf{沿屋面法线的一维运动。} 设有向法向坐标 \[
z=y\cos\alpha+x\sin\alpha.
\] 屋面满足 \(z=h\cos\alpha\)，抛体初速度法向分量 \[
v_{n0}=v_0\sin(\theta+\alpha),
\] 重力法向分量为 \(-g\cos\alpha\)。所以 \[
z(t)=v_0\sin(\theta+\alpha)t-\frac12g\cos\alpha\,t^2.
\] 这直接解释了最短时间为何要求
\(\theta+\alpha=\pi/2\)。在存在首次交点的范围，较小正根对 \(v_{n0}\)
单调递减。

\textbf{存在与相切。} 二次方程判别式 \[
\Delta=v_{n0}^2-2gh\cos^2\alpha.
\] \(\Delta>0\) 有两个形式交点，首次碰撞选小根；\(\Delta=0\)
是轨迹与直线相切；\(\Delta<0\)
根本达不到无限延长的屋面。屋面若只在有限区间，还须检查 \(x(t)\)
位于该区间，否则直线相交不等于实际碰撞。

最短时间 \[
t_{\min}=\frac{v_0-\sqrt{v_0^2-2gh\cos^2\alpha}}{g\cos\alpha}.
\] 为避免大速度下两近数相减，用等价形式 \[
\boxed{t_{\min}=
\frac{2h\cos\alpha}{v_0+\sqrt{v_0^2-2gh\cos^2\alpha}}}.
\] \(v_0^2\gg gh\) 时 \[
t_{\min}=\frac{h\cos\alpha}{v_0}
\left[1+\frac{gh\cos^2\alpha}{2v_0^2}+O(v_0^{-4})\right].
\] 第一项是忽略重力的法向飞行时间，第二项是向下重力延迟上升碰撞的修正。

\textbf{用交点坐标独立检验速度。} 首次交点 \(P\) 已知时， \[
v_x=v_0\cos\theta,\qquad
v_y=v_0\sin\theta-gt_P.
\] 速度大小的能量检查 \[
v_P^2=v_0^2-2gy_P
\] 以发射点为势能零点；\(h\) 是屋面截距，不是质点初始高度。把 \(h-y_P\)
填入势能差会多出一项 \(2gh\)。

法向碰撞速度 \[
v_n(t_P)=v_{n0}-g\cos\alpha\,t_P=\sqrt{\Delta}
\] 为正，说明首次交点仍向屋面上方运动；第二根给
\(-\sqrt{\Delta}\)，代表已穿过无限直线后的返回交点，实际首次撞击问题不会到达它。

\textbf{另一种优化条件。}
若目标改为``碰撞处速度垂直屋面''，不仅要求初始速度沿法线。屋面切向速度为
\[
v_\parallel(t)=v_0\cos(\theta+\alpha)+g\sin\alpha\,t.
\] 应同时满足 \(v_\parallel(t_P)=0\)
和交点方程。重力改变切向速度，所以它通常与最短时间角不同。图2-2的速度箭头应按实际
\(v_x,v_y\) 作切线，不能为了显示最优方向把碰撞速度画成法线。

\textbf{撞击后的速度：加入恢复系数。}
这是把屋面视为固定、光滑平面的变式。取向上法线和沿屋面向右下的单位切向量
\[
\boldsymbol n=(\sin\alpha,\cos\alpha),\qquad
\boldsymbol t=(\cos\alpha,-\sin\alpha).
\] 撞前速度分解为 \[
\boldsymbol v^-=v_n^-\boldsymbol n+v_t^-\boldsymbol t,
\] \[
v_n^-=v_0\sin(\theta+\alpha)-gt_P\cos\alpha,\qquad
v_t^-=v_0\cos(\theta+\alpha)+gt_P\sin\alpha.
\] 第一次从下方撞到屋面时 \(v_n^->0\)。若法向恢复系数为
\(e\)，光滑面不改变切向动量： \[
v_n^+=-e v_n^-,\qquad v_t^+=v_t^-.
\] 所以 \[
\boxed{\boldsymbol v^+=
\boldsymbol v^--(1+e)(\boldsymbol v^-\!\cdot\boldsymbol n)\boldsymbol n}.
\] 沿坐标轴展开得 \[
v_x^+=v_x^--(1+e)v_n^-\sin\alpha,
\] \[
v_y^+=v_y^--(1+e)v_n^-\cos\alpha.
\]
这两式可直接用图2-2的撞前切线检验；反弹方向一般既不竖直，也不与屋面垂直。

\textbf{耗能与相切极限。} 撞击瞬间位置不变，重力势能不变，动能损失为 \[
K^--K^+=\frac m2(1-e^2)(v_n^-)^2.
\] 又有 \(v_n^-=\sqrt{\Delta}\)，所以 \[
\boxed{\Delta K=\frac m2(1-e^2)\Delta}.
\] \(e=1\) 为弹性反射，速率不变；\(e=0\)
只消除法向速度，切向速度仍在。相切时
\(\Delta=0\)，法向冲量为零，因此``碰到直线''不一定发生有限冲量的撞击。

\textbf{反弹后会不会再次撞到同一无限平面。}
以第一次撞点为新原点，令撞后经过时间为 \(\tau\)。法向位移为 \[
z_{\rm rel}(\tau)=-e v_n^-\tau-\frac12g\cos\alpha\,\tau^2.
\] 对 \(0\le\alpha<\pi/2\)、\(e\ge0\)，当 \(\tau>0\)
时该位移非正；除退化的零速度情形外严格为负。因而在仅有重力的模型里，反弹后不会再次撞到这条屋面直线。若计算得到另一个正碰撞时间，通常是仍在使用撞前速度或误选了法线方向。

反弹后的最高点若存在，要满足 \(v_y^+>0\)，其相对第一次撞点的上升量为 \[
\Delta y_{\max}=\frac{(v_y^+)^2}{2g}.
\] 若
\(v_y^+\le0\)，则撞后最高点就是撞点本身，不能把平方公式当作实际未来上升高度。

\needspace{14\baselineskip}

\textbf{恰好可达屋面的临界发射。} 让初速度大小也可调，最小可达速度为 \[
v_{\rm crit}=\sqrt{2gh}\cos\alpha.
\] 发射方向取法线 \(\theta=\pi/2-\alpha\)，法向最高点恰好落在屋面。此时
\[
t_{\rm crit}=\frac{v_{\rm crit}}{g\cos\alpha}
=\sqrt{\frac{2h}{g}}.
\] 交点坐标为 \[
x_{\rm crit}=h\sin2\alpha,\qquad
y_{\rm crit}=h\cos2\alpha.
\] 代入 \(y=h-x\tan\alpha\) 可检验几何。碰撞前的速率为 \[
v_P^2=v_{\rm crit}^2-2gy_{\rm crit}
=2gh\sin^2\alpha.
\]
它通常不为零，而法向速度为零，因此这是一条沿屋面切向运动的相切轨迹。只有水平顶棚
\(\alpha=0\) 的特殊情形，交点同时是抛体的静止最高点。

\newpage

\section{第 3 题}\label{ux7b2c-3-ux9898}

一副近视游泳眼镜使用平凹薄透镜，折射率为
\(n\)，朝向水的一面为平面，另一面曲率半径为有符号量 \(R\)。

（1）求眼镜度数与 \(R\) 的关系；

（2）若佩戴者在水下看远处物体清楚，回到空气中继续佩戴时，远处物体是否仍清楚？

\begin{figure}[H]
\centering
\includegraphics[width=0.8\linewidth,height=\textheight,keepaspectratio,alt={图3-1　平凹镜片的平面入射面与曲面出射面}]{../figures/question-03.pdf}
\caption*{图3-1　平凹镜片的平面入射面与曲面出射面}
\end{figure}

\subsection{一、题型定位}\label{ux4e00ux9898ux578bux5b9aux4f4d-2}

本题考查单球面折射与平面界面的作用。水侧表面为平面，来自无穷远的平行光穿过该表面后仍保持平行；全部屈光能力由玻璃到眼侧空气的曲面产生，因此远物焦度不受水侧外界介质改变。

\subsection{二、知识点}\label{ux4e8cux77e5ux8bc6ux70b9-2}

\begin{enumerate}
\def\labelenumi{\arabic{enumi}.}
\tightlist
\item
  球面折射公式；
\item
  平面界面不改变平行光束会聚度；
\item
  有符号曲率半径；
\item
  光焦度与眼镜度数；
\item
  远物与有限物距的区别。
\end{enumerate}

\subsection{三、物理图像}\label{ux4e09ux7269ux7406ux56feux50cf-2}

平行光先从水或空气垂直通过平面表面，仍为平行光；到达凹球面后从折射率
\(n\) 的材料射向空气，形成虚焦点，起近视镜所需的发散作用。

\subsection{四、解题入口}\label{ux56dbux89e3ux9898ux5165ux53e3-2}

球面折射的近轴公式可写为 \[
\frac{n_2}{v}-\frac{n_1}{u}
=\frac{n_2-n_1}{R}. \tag{1}
\] 进入曲面前光束平行，故 \(u=\infty\)；此处 \(n_1=n,n_2=1\)。

\subsection{五、解析}\label{ux4e94ux89e3ux6790-2}

代入式（1）： \[
\frac1f=\frac{1-n}{R}. \tag{2}
\] 光焦度 \[
\Phi=\frac1f=\frac{1-n}{R}. \tag{3}
\] 若中国眼镜``度数''记为每屈光度100度，则 \[
\boxed{\mathcal D=100\frac{1-n}{R}}, \tag{4}
\] 其中 \(R\) 用米并按约定取符号。平凹近视镜的 \(\Phi<0\)。

\begin{figure}[H]
\centering
\includegraphics[width=0.8\linewidth,height=\textheight,keepaspectratio,alt={图3-2　水中与空气中平行光的相同远物焦度}]{../figures/solution-03.pdf}
\caption*{图3-2　水中与空气中平行光的相同远物焦度}
\end{figure}

第二问中，远处物体入射到平面表面的是平行光。不论外侧是水还是空气，平面界面只改变单条斜光线的方向比例，不会使平行束获得会聚或发散的光焦度；进入镜片后仍为平行束。曲面一侧介质与曲率未变，所以式（2）的焦度不变。

因此在题设只比较远处清晰度、且眼镜与眼睛的相对位置不变时， \[
\boxed{\text{回到空气中仍可看清远处物体}}. \tag{5}
\]
若物体为有限距离，平面界面会改变物距的光学共轭关系，此结论不再直接成立。

\textbf{用逐面光焦度核对符号。} 平面曲率
\(R_1=\infty\)，焦度为零；曲面从玻璃到空气的贡献 \[
\Phi_2=\frac{n_{\rm out}-n_{\rm glass}}R
=\frac{1-n}{R}.
\] 所以总焦度就是式（3）。近视用发散镜必须有
\(\Phi<0\)，该物理要求可反查 \(R\) 的有向符号。

远物在平面水---玻璃界面前后都对应平行束，平面不产生有限会聚度，因此水中、空气中的远物焦度相同。有限物距时入射束有曲率，平面折射会改变该曲率，结论才失效。

\textbf{用近轴折射矩阵核验。} 使用光线状态 \((y,n\vartheta)\)。从
\(n_1\) 到 \(n_2\) 的球面折射矩阵 \[
S(\Phi)=\begin{pmatrix}1&0\\-\Phi&1\end{pmatrix},
\qquad \Phi=(n_2-n_1)/R.
\] 传播厚度 \(d\)、介质折射率 \(n\)： \[
T(d)=\begin{pmatrix}1&d/n\\0&1\end{pmatrix}.
\]
本题水侧平面曲率为零，故矩阵为单位阵；平行入射的斜率也为零。出射凹面球心在右，\(R>0,n_2=1<n\)，所以
\[
n_{\rm out}\vartheta_{\rm out}=-\Phi y
=\frac{n-1}{R}y.
\]
上方光线向上偏折、下方光线向下偏折，因此出射发散，虚焦点在左。图3-2反向延长线必须与此符号一致。

\textbf{水的作用与不能推广的结论。}
水侧平面不改变光线的约化倾角，平行入射时内部仍平行，所以后焦距不含
\(n_w\)： \[
f'=\frac{n_{\rm out}}{\Phi}=\frac{R}{1-n}.
\] 前焦距按不同入、出射介质定义为 \(f=-n_w/\Phi\)，它含
\(n_w\)。因此``不含水折射率''只适用于本题的后焦距，不能泛称透镜全部成像性质与水无关。

\textbf{有限物距。} 若先取薄元件近似，物点在水侧距离 \(s>0\)，像距
\(s'\) 向空气侧为正，成像式 \[
\frac{n_w}{s}+\frac1{s'}=\Phi.
\] 由于 \(\Phi<0\)，实物给 \(1/s'<0\)，像仍在左侧。横向放大率 \[
m=-\frac{n_w s'}s
\] 为正，故虚像正立。取 \(s\to\infty\) 回到
\(s'=1/\Phi\)，可核验符号口径。

\textbf{反转透镜时。} 若改成曲面先入水、平面出空气，第一面功率为
\((n-n_w)/R_1\)，并非原来的
\((1-n)/R\)。透镜两侧介质不同，翻转改变实际界面顺序，不能使用空气中薄透镜翻面焦距不变的经验。

\textbf{单位辨认。} 光学功率定义为 \(\Phi\)，单位
\(\mathrm{m^{-1}}\)。若题给焦距以厘米表示，则 \[
d=\frac{100}{f'_{\rm cm}}=\Phi_{\rm m^{-1}},\qquad
D=100d=\frac{10000}{f'_{\rm cm}}.
\] 这里 \(d\) 的单位是屈光度，\(D\) 是眼镜口语中的``度数''。若 \(R\)
用米，\(d=(1-n)/R\)，眼镜度数仍为
\(D=100(1-n)/R\)。两种计量相差一百倍，数值代入前应同时写清长度单位和所求计量。

\textbf{保留玻璃厚度的成像变式。}
原题平行入射的后焦距不受玻璃厚度影响。若物体距离有限，这个结论不能照搬。设平面到曲面顶点的轴向厚度为
\(t\)，仍用光线坐标 \((y,u)\)，其中 \(u=n\vartheta\)。元件矩阵为 \[
M=S(\Phi)T(t)=
\begin{pmatrix}
1&t/n\\
-\Phi&1-\Phi t/n
\end{pmatrix}.
\]
平面界面没有曲率，约化倾角连续，因此不另加功率项。若轴上物点距第一平面
\(s>0\)，一条高度为 \(y\) 的光线进入元件时 \[
u_{\rm in}=\frac{n_w}{s}y.
\] 抵达最后一个面的高度和出射约化倾角为 \[
y_{\rm out}=\left(1+\frac{tn_w}{ns}\right)y,
\] \[
u_{\rm out}=\left[-\Phi+
\left(1-\frac{\Phi t}{n}\right)\frac{n_w}s\right]y.
\] 空气中像点满足 \(y_{\rm out}+s'u_{\rm out}=0\)，故 \[
\boxed{s'=
-\frac{s/n_w+t/n}{1-\Phi(s/n_w+t/n)}}.
\] 取 \(t\to0\)，还原薄元件公式；取 \(s\to\infty\)，得到
\(s'\to1/\Phi\)。因为 \(\Phi<0\) 且
\(s,t>0\)，分母为正，像距始终为负，正是图3-2所示的虚像支路。

\textbf{横向放大率也要算传播。} 对离轴物点
\(y_o\)，第一面的入射约化倾角满足 \[
u_{\rm in}=\frac{n_w}{s}(y-y_o).
\] 总成像矩阵为 \[
T_{\rm air}(s')\,M\,T_{\rm water}(s).
\] 在像平面，总矩阵的右上角为零，左上角就是横向放大率： \[
\boxed{\mathfrak m=
\frac1{1-\Phi(s/n_w+t/n)}}.
\] \(\Phi<0\) 时 \(0<\mathfrak m<1\)，虚像正立且缩小。薄元件极限下 \[
\mathfrak m=-\frac{n_ws'}s,
\]
与前面的像距式一致；有厚度时应使用完整矩阵式，不能直接把第一面物距与最后一面像距代入薄元件放大率公式。

\textbf{一次带单位的代入。} 取 \(n=1.50,R=0.10\,\mathrm m\)，则 \[
\Phi=-5.00\,\mathrm{m^{-1}},\qquad f'=-0.20\,\mathrm m.
\] 这是 \(-5.00\) 屈光度、通常所说的 \(-500\) 度。再取
\(n_w=4/3,s=0.30\,\mathrm m,t=0.01\,\mathrm m\)： \[
s/n_w+t/n=0.231667\,\mathrm m,
\] \[
s'=-0.107336\,\mathrm m,\qquad \mathfrak m=0.46332.
\] 这个例子同时检验介质折射率、长度单位和虚像符号。

\textbf{有限孔径的单球面光路。} 近轴模型用角度的一阶近似。对曲率半径
\(R\) 的出射凹面，设顶点为 \(x=0\)、球心为 \((R,0)\)，高度 \(y>0\)
的平行光到达球面坐标 \[
x_P=R-\sqrt{R^2-y^2}.
\] 局部法线相对光轴为 \(-\arcsin(y/R)\)，玻璃内入射角为
\(i=\arcsin(y/R)\)。精确出射角为 \[
\vartheta_{\rm out}
=\arcsin(ny/R)-\arcsin(y/R).
\] 出射反向延长线与光轴交点是 \[
\boxed{x_F(y)=x_P-\frac{y}{\tan\vartheta_{\rm out}}}.
\] \(y/R\to0\) 时 \(x_P\to0,\vartheta_{\rm out}\simeq(n-1)y/R\)，因此
\(x_F\to-R/(n-1)\)，回到近轴虚焦距。有限高度下 \(x_F\) 随 \(y\)
变化，图3-2的共同虚焦点只表示近轴阶。

还需满足
\(ny/R\le1\)。若超出这个条件，玻璃到空气界面发生全反射，不能继续用实数出射角。这个孔径限制与虚像的正负号是两个独立问题。

\needspace{14\baselineskip}

\textbf{同一片镜片，远物焦度相同，近物成像可以不同。} 取补充参数
\(n=1.50\)、\(R=0.25\,\mathrm m\)、厚度
\(t=3.0\,\mathrm{mm}\)，曲面后为空气，则
\(\Phi=(1-n)/R=-2.0\,\mathrm{m^{-1}}\)。物点在平面前
\(s=0.50\,\mathrm m\)，入射介质折射率为
\(n_w\)。从物平面到曲面后的约化传播长度为 \[
B_0=\frac{s}{n_w}+\frac{t}{n}.
\] 再在空气中传播有向距离 \(v\)，成像要求总矩阵的右上元为零： \[
B_0+v(1-\Phi B_0)=0.
\] 所以 \[
v=-\frac{B_0}{1-\Phi B_0},\qquad
m_y=\frac1{1-\Phi B_0}.
\] 在空气中 \(n_w=1\) 时，得到
\(v\simeq-0.25050\,\mathrm m\)、\(m_y\simeq0.49900\)；水侧取
\(n_w=1.33\) 时，得到
\(v\simeq-0.21524\,\mathrm m\)、\(m_y\simeq0.56952\)。两者均为正立缩小的虚像，但像距不同。

当
\(s\to\infty\)，\(v\to1/\Phi=-0.50\,\mathrm m\)，两种入射介质才回到相同的后焦点。由此可见，原题的结论依赖平行入射；对有限物距，平面界面虽然不提供焦度，仍改变传播矩阵中的角度与距离关系。

\newpage

\section{第 4 题}\label{ux7b2c-4-ux9898}

（1）Gilbert 模型由相距 \(d\) 的磁荷 \(\pm q_m\) 构成，磁偶极矩
\(p_m=q_md\)。求轴线上 \(z\gg d\) 的场。 （2）Ampere 模型为半径
\(r\)、电流 \(I\) 的圆环，\(p_m=I\pi r^2\)。求轴线上 \(z\gg r\) 的场。
（3）半径 \(R\)、长度 \(L\) 的圆柱中，单位体积有 \(n\) 个沿 \(+z\)
的同向偶极子，满足
\(R\gg d,L\gg d\)。分别求两种微观模型在圆柱中心的场。原题未给
\(L\gg R\)，不可擅加长螺线管条件。

\includegraphics[width=0.86\linewidth,height=\textheight,keepaspectratio,alt={题4　题设与符号}]{../figures/question-04.pdf}

\subsection{一、题型定位}\label{ux4e00ux9898ux578bux5b9aux4f4d-3}

两种偶极源的远场与内部场比较。相同偶极矩使远场一致，圆柱内部的局部源结构却不同。把假想真实磁荷模型与实际磁化中的辅助磁荷混用，会抹去题目要比较的差别。

\subsection{二、知识点}\label{ux4e8cux77e5ux8bc6ux70b9-3}

偶极远场；有限圆盘轴线场；有限螺线管；磁化强度；退磁场；磁感应强度与辅助磁场；长细与扁平极限。

\subsection{三、物理图像}\label{ux4e09ux7269ux7406ux56feux50cf-3}

图4-1磁矩沿 \(+z\)。Gilbert
圆柱顶面为正磁荷、底面为负磁荷，中心场从顶向底，沿 \(-z\)；Ampere
圆柱的侧面环电流给中心沿 \(+z\) 的场。图4-2保留同一
\(R,L\)，没有把端面画成无限大平面。

\subsection{四、解题入口}\label{ux56dbux89e3ux9898ux5165ux53e3-3}

先分别展开单偶极远场，再令 \(M=np_m\)。中心到每个端面距离
\(a=L/2\)，由有限圆盘积分求 Gilbert 场；把侧面电流拆为连续圆环求 Ampere
场。

\subsection{五、解析}\label{ux4e94ux89e3ux6790-3}

\textbf{对称磁荷的精确轴线场。} 把正、负磁荷放在 \(z=\pm d/2\)，取观察点
\(z>d/2\)： \[
B_G(z)=\frac{\mu_0q_m}{4\pi}
\left[\frac1{(z-d/2)^2}-\frac1{(z+d/2)^2}\right].
\] 通分得 \[
B_G(z)=\frac{\mu_0p_m}{2\pi z^3}
\left(1-\frac{d^2}{4z^2}\right)^{-2}.
\] 所以 \[
\boxed{B_G\simeq\frac{\mu_0p_m}{2\pi z^3}},
\] 方向沿磁矩。下一项相对修正为
\(+d^2/(2z^2)\)。坐标原点若取在负磁荷处，精确式的形状会变，但在真正远场领先项相同。

\textbf{圆电流的精确场。} \[
B_A(z)=\frac{\mu_0Ir^2}{2(z^2+r^2)^{3/2}}
=\frac{\mu_0p_m}{2\pi z^3}
\left(1+\frac{r^2}{z^2}\right)^{-3/2}.
\] 因此远场同样为 \[
\boxed{B_A\simeq\frac{\mu_0p_m}{2\pi z^3}}.
\] 下一项相对修正为
\(-3r^2/(2z^2)\)。两个远场的领先项相同，修正的符号与长度尺度不同，不能声称全部空间的场完全一致。

\textbf{磁化与端面源。} 密集同向偶极子的磁矩体密度为 \[
M=np_m.
\] 在小偶极尺寸相对 \(R,L\) 可忽略时，Gilbert
模型内部磁荷抵消，仅两端留下面磁荷 \(\pm M\)。以半径 \(s\)
的圆环微元积分，正盘在距盘面 \(a>0\) 的轴线点产生场大小 \[
B_{\rm disk}=\frac{\mu_0M}{2}
\left(1-\frac a{\sqrt{a^2+R^2}}\right).
\] 中心处两个端面均使场向 \(-z\)，所以 \[
\boxed{\boldsymbol B_G(0)=
-\mu_0M\left(1-\frac a{\sqrt{a^2+R^2}}\right)\boldsymbol e_z},
\qquad a=L/2.
\]
负号不是磁矩方向定义改变：顶面正磁荷的场在下方指向下，底面负磁荷的场在上方也指向下。

\begin{figure}[H]
\centering
\includegraphics[width=0.86\linewidth,height=\textheight,keepaspectratio,alt={图4-2　有限圆柱中心的两种场与局部差值}]{../figures/solution-04.pdf}
\caption*{图4-2　有限圆柱中心的两种场与局部差值}
\end{figure}

\textbf{侧面电流积分。} Ampere 微观环的内部边电流相消，只留下环向面电流
\(K_b=M\)。高为 \(dz\) 的侧面条带对应电流 \(dI=M\,dz\)。中心场 \[
B_A(0)=\int_{-a}^{a}
\frac{\mu_0MR^2\,dz}{2(R^2+z^2)^{3/2}}.
\] 利用原函数 \(z/[R^2\sqrt{R^2+z^2}]\)，得到 \[
\boxed{\boldsymbol B_A(0)=
\mu_0M\frac a{\sqrt{a^2+R^2}}\boldsymbol e_z}.
\] 两种字面微观源的中心场差为 \[
\boxed{\boldsymbol B_A-\boldsymbol B_G=\mu_0\boldsymbol M}.
\] 原稿把两者都写成 \(\mu_0M\)，既忽略有限长几何，也把磁荷场方向写反了。

\textbf{与实际磁化辅助场的关系。}
实际材料的磁化由电流源描述。用等效端面磁荷计算的是辅助场
\(\boldsymbol H_G\)，其中心值为 \[
\boldsymbol H_G=-M\left(1-\frac a{\sqrt{a^2+R^2}}\right)\boldsymbol e_z.
\] 随后应恢复 \[
\boldsymbol B=\mu_0(\boldsymbol H_G+\boldsymbol M)
=\boldsymbol B_A.
\]
因此``辅助磁荷法与束缚电流法得到同一实际磁场''成立；``假想真实磁荷模型与电流模型在材料内部的
\(\boldsymbol B\) 相同''不成立。两者差别在于是否已加入局部磁化项，可参看
\href{https://www.feynmanlectures.caltech.edu/II_36.html}{Feynman
Lectures II，第36章}。

\textbf{长细与薄盘极限。} 若另有 \(L\gg R\)， \[
\frac a{\sqrt{a^2+R^2}}
=1-\frac{R^2}{2a^2}+\cdots,
\] 于是 \(B_A\to\mu_0M\)，而 \(B_G\to0^-\)。反之 \(L\ll R\)，有 \[
B_A\simeq\mu_0M L/(2R),\qquad
B_G\simeq-\mu_0M.
\]
无限大端面近似适合薄盘而非长细圆柱；原稿把这两个相反极限混在了一起。题给
\(L\gg d\) 只说明偶极尺寸小，不提供圆柱长径比。

\textbf{任意轴线点的有限场。} 在 \(-a<z<a\) 的圆柱内部， \[
B_A(z)=\frac{\mu_0M}{2}
\left[
\frac{z+a}{\sqrt{R^2+(z+a)^2}}
-\frac{z-a}{\sqrt{R^2+(z-a)^2}}
\right].
\] Gilbert 字面磁荷场满足
\(B_G(z)=B_A(z)-\mu_0M\)。在圆柱外没有磁化项，两模型的连续偶极极限给相同外部场。这个关系比仅检查中心更强，能说明差别来自局部源，而不是中心计算的偶然误差。

\textbf{宏观总磁矩的远场。} 整个圆柱总偶极矩为 \[
p_{\rm tot}=M\pi R^2L.
\] 当 \(|z|\gg R,L\)，两模型都回到 \[
B_z\simeq\frac{\mu_0p_{\rm tot}}{2\pi|z|^3}
\] （在轴的两侧均沿 \(+z\)）。中心内部场不同与远处总偶极矩相同并不矛盾。

\textbf{退磁系数的轴向积分。} 令 \(a=L/2\)。端面磁荷密度
\(\pm M\)，面积元磁荷 \(M\,2\pi s\,ds\)。一个端面在中心的轴向场大小 \[
\frac{\mu_0M}{4\pi}\int_0^R
\frac{a\,2\pi s\,ds}{(a^2+s^2)^{3/2}}
=\frac{\mu_0M}{2}\left(1-\frac a{\sqrt{a^2+R^2}}\right).
\] 两个端面方向相同，沿
\(-z\)。这说明有限半径积分的上限不可替换成无穷，除非另给 \(R\gg a\)。

侧面等效电流 \(K=M\)，每一厚度 \(dz\) 的电流环为 \(dI=M\,dz\)。中心场 \[
B_A=\frac{\mu_0M}{2}\int_{-a}^{a}
\frac{R^2\,dz}{(R^2+z^2)^{3/2}}
=\mu_0M\frac a{\sqrt{a^2+R^2}}.
\] 两种源各自的积分核不同，却精确满足
\(B_A-B_G=\mu_0M\)。这不是只在长圆柱极限出现的关系。

\textbf{长细和扁平极限分开。} 令 \(x=a/R\)，则 \[
\frac{B_A}{\mu_0M}=\frac{x}{\sqrt{1+x^2}},
\qquad \frac{B_G}{\mu_0M}=\frac{x}{\sqrt{1+x^2}}-1.
\] 长细 \(x\gg1\)： \[
B_A=\mu_0M\left[1-\frac1{2x^2}+\frac3{8x^4}+\cdots\right],
\] \[
B_G=-\mu_0M\left[\frac1{2x^2}-\frac3{8x^4}+\cdots\right].
\] 扁平 \(x\ll1\)： \[
B_A=\mu_0M[x-\tfrac12x^3+\cdots],\qquad
B_G=-\mu_0M[1-x+\tfrac12x^3+\cdots].
\] ``偶极尺寸远小于 \(R,L\)''只控制连续化，无法决定 \(x\) 大小。

\textbf{中心不是全体内部。} 一般轴上位置 \(z\) 的 Ampere 场 \[
B_A(z)=\frac{\mu_0M}{2}
\left[\frac{z+a}{\sqrt{(z+a)^2+R^2}}
-\frac{z-a}{\sqrt{(z-a)^2+R^2}}\right].
\] 中心 \(z=0\) 回到正文。圆柱内任一点处，辅助磁荷计算
\(H\)，真实磁化场满足 \(B=\mu_0(H+M)\)；圆柱外
\(M=0\)，两模型的远场相同。不能把中心式视为整个有限圆柱内的均匀场。

\textbf{远场由总磁矩控制。} 连续化总磁矩 \[
p_{\rm tot}=M\pi R^2L.
\] 当 \(z\gg R,L\)，轴上场 \[
B(z)\simeq\frac{\mu_0p_{\rm tot}}{2\pi z^3}.
\] 由一般轴式展开得到相同首项，成为端面法与电流法的远场交叉检验。

\textbf{轴线上任意位置的场。} 把中心结果推广到
\(z\)，可检查磁荷与电流模型何处相同、何处相差。半长为 \(a\)、半径为
\(R\)，均匀磁化沿 \(+z\)，侧面的束缚面电流密度为 \(K_b=M\)。半径 \(R\)
的薄电流环在轴上的场为 \[
dB_A(z)=\frac{\mu_0M R^2\,dz'}{2[R^2+(z-z')^2]^{3/2}}.
\] 对 \(-a\le z'\le a\) 积分： \[
\boxed{B_A(z)=\frac{\mu_0M}{2}\left[
\frac{z+a}{\sqrt{(z+a)^2+R^2}}
-\frac{z-a}{\sqrt{(z-a)^2+R^2}}\right]}.
\] \(z=0\) 就是正文的中心结果。由于括号关于 \(z\)
为偶函数，两端外部轴向场同向，这符合轴向磁偶极场的方向。

\textbf{辅助磁荷给出的量。} 端面磁荷为
\(\sigma_m=\pm M\)。单个均匀圆盘在距离 \(d\) 处的轴向辅助场，按 \(d\)
的正负写成 \[
H_{\rm disk}(d)=\frac{\sigma_m}{2}
\left[\operatorname{sgn}d-\frac d{\sqrt{d^2+R^2}}\right].
\] 两端相加。在内部 \(-a<z<a\)： \[
H_G(z)=\frac M2\left[
\frac{z+a}{\sqrt{(z+a)^2+R^2}}
-\frac{z-a}{\sqrt{(z-a)^2+R^2}}-2\right].
\] 所以内部有 \[
B_A(z)=\mu_0[H_G(z)+M].
\] 在外部 \(|z|>a\)，两个符号项相消，故 \[
B_A(z)=\mu_0H_G(z).
\]
这说明外场等价并不推出内部两个量相同。图4-2中的箭头比较只针对内部中心；若把同样的反向箭头画到远处，会违背外场等价。

\textbf{远场的偶极检验。} 对 \(z\gg a,R\) 展开轴向公式，首个不消失的项为
\[
B_A(z)=\frac{\mu_0MR^2a}{z^3}+O(z^{-5}).
\] 磁体的总磁矩为 \[
\mathcal m=M(2a\pi R^2).
\] 轴向磁偶极公式给 \[
B_{\rm dip}=\frac{\mu_0}{4\pi}\frac{2\mathcal m}{z^3}
=\frac{\mu_0MR^2a}{z^3},
\] 与积分一致。远场若出现 \(z^{-2}\)
项，就说明两个端面的磁荷没有等量异号，或者积分漏了一端。

\needspace{14\baselineskip}

\textbf{两个内部场何时等大。} 中心场满足 \[
\frac{B_A}{\mu_0M}=\frac{a}{\sqrt{a^2+R^2}},
\qquad
\frac{|B_G|}{\mu_0M}=1-\frac{a}{\sqrt{a^2+R^2}}.
\] 令两大小相等，可得 \[
\boxed{a/R=1/\sqrt3}.
\] 长柱 \(a\gg R\) 时 \[
B_A=\mu_0M\left[1-\frac{R^2}{2a^2}+O(R^4/a^4)\right],
\] \[
B_G=-\mu_0M\left[\frac{R^2}{2a^2}+O(R^4/a^4)\right].
\] 薄片 \(a\ll R\) 时则为 \[
B_A\simeq\mu_0M\,a/R,\qquad
B_G\simeq-\mu_0M(1-a/R).
\] 两种极限都保持
\(B_A-B_G=\mu_0M\)，但主导的场不同。图中的箭头方向应固定，长度可随长径比改变，不能仅以``有限磁体''给出一个与几何无关的同等长示意。

\newpage

\section{第 5 题}\label{ux7b2c-5-ux9898}

如图所示，刚性密闭容器内有 \(1\,\mathrm{mol}\) 液态水和
\(1\,\mathrm{mol}\)
水蒸气，两相处于平衡状态。容器内装有电热丝，器壁绝热。初态温度为
\(T_0=100^\circ\mathrm C\)。已知水的摩尔汽化热
\(L=40.8\,\mathrm{kJ\,mol^{-1}}\)，液态水的摩尔热容
\(C=75.8\,\mathrm{J\,K^{-1}\,mol^{-1}}\)，水蒸气的定容摩尔热容
\(C_v=3R\)，普适气体常量
\(R=8.31\,\mathrm{J\,K^{-1}\,mol^{-1}}\)。现用电热丝加热，使系统温度升高
\(1^\circ\mathrm C\)。假定 \(L\)
的变化、液态水体积均可忽略，水蒸气可视为理想气体。以下结果保留三位有效数字。

（1）已知饱和蒸气压满足 \(dp/dT=L/(TV_m)\)，其中 \(V_m\)
为水蒸气的摩尔体积，求加热后水蒸气摩尔数的增量 \(\Delta n\)。

（2）求电热丝提供的热量 \(Q\)。

\begin{figure}[H]
\centering
\includegraphics[width=0.462\linewidth,height=\textheight,keepaspectratio,alt={题5　刚性容器中的液---汽平衡系统}]{../figures/question-05.pdf}
\caption{题5　刚性容器中的液---汽平衡系统}
\end{figure}

\subsection{一、题型定位}\label{ux4e00ux9898ux578bux5b9aux4f4d-4}

升温同时改变两相的温度和物质的分配。系统始终处在液态水与饱和水蒸气共存状态，升温会同时引起两类效应：原有物质的温度升高，以及一部分液态水继续汽化。第一问利用克劳修斯---克拉珀龙关系求汽化量；第二问必须在刚性、绝热容器的条件下正确区分焓变与内能变。

\subsection{二、知识点}\label{ux4e8cux77e5ux8bc6ux70b9-4}

\begin{enumerate}
\def\labelenumi{\arabic{enumi}.}
\tightlist
\item
  饱和蒸气压随温度的变化；
\item
  理想气体状态方程及其微分；
\item
  摩尔汽化热 \(L\) 的热力学含义；
\item
  刚性密闭系统的热力学第一定律；
\item
  小温差条件下的一阶近似。
\end{enumerate}

\subsection{三、物理图像}\label{ux4e09ux7269ux7406ux56feux50cf-4}

容器中原有 \(1\,\mathrm{mol}\) 液态水和 \(1\,\mathrm{mol}\)
饱和水蒸气。电热丝向系统输入能量，容器壁绝热且总体积不变；升温
\(1\,\mathrm{K}\)
后，饱和蒸气压升高，必须有少量液体汽化才能达到新的相平衡。

\subsection{四、解题入口}\label{ux56dbux89e3ux9898ux5165ux53e3-4}

先把蒸气摩尔数写成 \[
n=\frac{pV}{RT},
\] 再对温度求导，并用题给的饱和蒸气压关系消去
\(dp/dT\)。求热量时以容器内全部水为系统；因为容器刚性，外界体积功为零，所以电热丝提供的能量等于系统内能的增加。

\subsection{五、解析}\label{ux4e94ux89e3ux6790-4}

已知初温 \[
T_0=100^\circ\mathrm C=373\,\mathrm K,
\] 升温量为 \[
\Delta T=1\,\mathrm K.
\]

\needspace{3\baselineskip}
\textbf{（1）水蒸气摩尔数的增量}

忽略液态水体积，容器体积等于水蒸气所占体积。对理想水蒸气有 \[
pV=nRT. \tag{1}
\] 容器刚性，所以 \(V\) 为常量。将式（1）写成 \[
n=\frac{pV}{RT}. \tag{2}
\] 对 \(T\) 求导： \[
\begin{aligned}
\frac{dn}{dT}
&=\frac{V}{R}\frac{d}{dT}\left(\frac pT\right)\\
&=\frac{V}{R}\left(\frac1T\frac{dp}{dT}-\frac p{T^2}\right). \end{aligned} \tag{3}
\]

题给液---汽平衡关系 \[
\frac{dp}{dT}=\frac{L}{TV_m}, \tag{4}
\] 其中理想蒸气的摩尔体积 \[
V_m=\frac{RT}{p}. \tag{5}
\] 把式（5）代入式（4）： \[
\frac{dp}{dT}
=\frac{L}{T}\frac p{RT}
=\frac{Lp}{RT^2}. \tag{6}
\] 再将式（6）代入式（3）： \[
\begin{aligned}
\frac{dn}{dT}
&=\frac{V}{R}\left(\frac{Lp}{RT^3}-\frac p{T^2}\right)\\
&=\frac{pV}{RT}\frac1T\left(\frac{L}{RT}-1\right)\\
&=\frac nT\left(\frac{L}{RT}-1\right). \end{aligned} \tag{7}
\] 这里第二行使用了 \(pV=nRT\)。

由于 \(\Delta T/T_0\ll1\)，只保留关于 \(\Delta T\)
的一阶量，并在初态取各系数。初态蒸气摩尔数 \(n_0=1\,\mathrm{mol}\)，故
\[
\Delta n\simeq
\left.\frac{dn}{dT}\right|_{T_0}\Delta T
=\frac{n_0\Delta T}{T_0}
\left(\frac{L}{RT_0}-1\right). \tag{8}
\] 代入数据： \[
\begin{aligned}
\frac{L}{RT_0}
&=\frac{40.8\times10^3}{8.31\times373}\\
&\approx13.16,\\[2mm]
\Delta n
&=\frac{1}{373}(13.16-1)\,\mathrm{mol}\\
&\approx3.26\times10^{-2}\,\mathrm{mol}.
\end{aligned}
\] 所以 \[
\boxed{\Delta n=0.0326\,\mathrm{mol}}. \tag{9}
\]

\begin{figure}[H]
\centering
\includegraphics[width=0.8\linewidth,height=\textheight,keepaspectratio,alt={图5-2　微小升温过程中能量的两部分去向}]{../figures/solution-05.pdf}
\caption*{图5-2　微小升温过程中能量的两部分去向}
\end{figure}

\needspace{3\baselineskip}
\textbf{（2）电热丝提供的热量}

电热丝的供能分成两部分：

\begin{enumerate}
\def\labelenumi{\arabic{enumi}.}
\tightlist
\item
  初态的 \(1\,\mathrm{mol}\) 液态水和 \(1\,\mathrm{mol}\)
  水蒸气整体升温；
\item
  另有 \(\Delta n\) 摩尔液态水在近似温度 \(T_0\) 下转变成蒸气。
\end{enumerate}

第一部分引起的内能增量为 \[
\Delta U_{\rm sens}
=\bigl(C+3R\bigr)\Delta T. \tag{10}
\]

下面确定每摩尔汽化所需的内能。摩尔汽化热 \(L\) 是定压汽化的焓变： \[
L=\Delta h=\Delta u+p\Delta V_m. \tag{11}
\] 忽略液态水摩尔体积，\(\Delta V_m\simeq V_m\)；理想蒸气满足
\(pV_m=RT_0\)，所以 \[
\Delta u_{\rm vap}=L-RT_0. \tag{12}
\]
这一步是本题的关键：容器刚性，系统不向外界做总体积功，能量守恒中应使用汽化内能
\(L-RT_0\)，不能直接把 \(L\) 当成系统内能增量。

于是 \[
Q=\Delta U
=\bigl(C+3R\bigr)\Delta T
+\Delta n\bigl(L-RT_0\bigr). \tag{13}
\] 逐项计算： \[
\begin{aligned}
C+3R
&=75.8+3\times8.31\\
&=100.73\,\mathrm{J\,K^{-1}},\\
L-RT_0
&=40.8\times10^3-8.31\times373\\
&\approx3.770\times10^4\,\mathrm{J\,mol^{-1}},\\
Q
&=100.73\times1
+0.0326\times3.770\times10^4\\
&\approx1.33\times10^3\,\mathrm J.
\end{aligned}
\] 因此 \[
\boxed{Q=1.33\,\mathrm{kJ}}. \tag{14}
\]

量纲检查：式（8）中括号内无量纲，\(n_0\Delta T/T_0\)
的量纲为摩尔；式（13）两项均为能量。数量级上，汽化约
\(0.033\,\mathrm{mol}\)，所需能量约
\(1.2\,\mathrm{kJ}\)，远大于整体升温所需的约
\(0.10\,\mathrm{kJ}\)，结果合理。

\textbf{线性化的自洽性检查。} 得到 \(\Delta n\simeq0.0326\,{\rm mol}\)
后，要与初始1 mol液体比较： \[
\frac{\Delta n}{1\,{\rm mol}}\simeq3.3\%\ll1.
\] 因此液相尚未耗尽，两相平衡假设仍可使用；温升
\(\Delta T/T_0\simeq2.7\times10^{-3}\) 也足够小。若计算出 \(\Delta n\)
接近1 mol，就不能继续使用定系数的一阶展开。

总热量可分成整体升温和相变两部分： \[
Q_{\rm warm}=(C+3R)\Delta T,
\qquad
Q_{\rm vap}=(L-RT_0)\Delta n.
\] 分别约为 \(0.10\,{\rm kJ}\) 和 \(1.23\,{\rm kJ}\)，相加为
\(1.33\,{\rm kJ}\)。分项计算既说明汽化潜热主导，也能快速发现把 \(L\)
的千焦单位误当成焦耳的数量级错误。

\textbf{有限温升必须协调潜热与热容。}
本题只求一开尔文温升的一阶结果，使用初态 \(L_0\)
与热容足够。若进一步求有限温升，既令 \(C_l,C_v\) 恒定，又令 \(L\)
严格恒定会不相容。理想蒸气、液体体积忽略时 Kirchhoff 关系 \[
\frac{dL}{dT}=C_v+R-C_l=D.
\] 所以一致的延伸模型取 \[
L(T)=L_0+D(T-T_0).
\] 题给 \(C_v=3R\)，于是 \(D=4R-C_l\)。蒸气与液体摩尔内能差 \[
\Delta u(T)=L(T)-RT,
\] 其温度导数为 \(C_v-C_l\)，与两相热容相符。

\textbf{饱和压强的有限积分。} 用 \[
\frac{d\ln p}{dT}=\frac{L(T)}{RT^2},
\] 分离恒定与线性项： \[
\ln\frac{p(T)}{p_0}=
\frac{L_0-DT_0}{R}\left(\frac1{T_0}-\frac1T\right)
+\frac D R\ln\frac T{T_0}.
\] 刚性容器 \(n=pV/(RT)\)，所以 \[
\frac{n(T)}{n_0}=\frac{T_0}{T}\frac{p(T)}{p_0}.
\] 该式仅在两相同时存在时有效；当
\(n(T)=2\,\mathrm{mol}\)，液体耗尽，后续不再服从饱和压强。

\textbf{有限内能差。} 总物质的量为
\(n_{\rm tot}=2\)，选择液体初态摩尔内能为零： \[
U(T)=n_{\rm tot}C_l(T-T_0)+n(T)\Delta u(T).
\] 初态 \(U_0=n_0\Delta u(T_0)\)，所以电热丝输入 \[
Q=n_{\rm tot}C_l(T-T_0)
+n(T)\Delta u(T)-n_0\Delta u(T_0).
\] 对初态求导： \[
Q'(T_0)=C_l+C_v+
\left.\frac{dn}{dT}\right|_0(L_0-RT_0),
\] 恰为正文的一阶答案。这里
\(2C_l+n_0(C_v-C_l)=C_l+C_v\)，避免重复计入新汽化部分的升温。

\textbf{蒸气增加率的温度变化。} 令 \[
a(T)=\frac1T\left[\frac{L(T)}{RT}-1\right].
\] 则 \(n'=na\)， \[
n''=n(a^2+a'),\quad
a'=\frac{D}{RT^2}-\frac{2L}{RT^3}+\frac1{T^2}.
\] 因此二阶温升修正为
\(n_0(a_0^2+a'_0)(\Delta T)^2/2\)。能否保留三位有效数字应比较这个修正与首阶项，而非只看
\(\Delta T/T_0\)。

\textbf{物理热量口径。}
容器外边界刚性，外部总体积功为零；局部汽化仍需推动蒸气，因此相变焓 \(L\)
与系统内能增量差
\(RT\)。若题目改成恒压活塞容器，第一定律还需加活塞功，热量才在适当条件下包含相变焓本身。

\textbf{恒温改变容器体积。}
这是与原题``固定体积升温''互补的变式。设液、汽两相共存，液体体积忽略，恒温
\(T\) 时饱和压强不变： \[
p=p_s(T),\qquad n_v=\frac{p_sV}{RT}.
\] 因此 \[
dn_v=\frac{p_s}{RT}\,dV.
\] 汽化焓为 \(L\)，一摩尔液体变成汽体时的内能增量为 \(L-RT\)，所以 \[
dU=(L-RT)\,dn_v.
\] 对外准静态作功是 \(p_s\,dV=RT\,dn_v\)。第一定律给 \[
\boxed{\delta Q=dU+p_s\,dV=L\,dn_v
=\frac{Lp_s}{RT}\,dV}.
\] 原题固定体积没有外界体积功，潜热项应使用
\(L-RT\)；这个恒温变式存在体积功，热量才是
\(L\,dn_v\)。两种过程用不同热量表达式并不矛盾。

\textbf{两相存在的体积区间。} 若总物质的量为 \(N\)，两相共存必须有 \[
0<n_v<N.
\] 所以在忽略液体体积的模型里，上界是 \[
V_{\rm dry}=\frac{NRT}{p_s(T)}.
\] 当 \(V\) 增大到此值，液体恰好全部汽化。继续增大体积后进入单一汽相：
\[
p=\frac{NRT}{V}<p_s(T).
\] 恒温 \(p\)---\(V\)
图先是平台，再接下降双曲线。平台结束处压强连续，但斜率改变；不能把两相平台延伸到任意大的体积。

\textbf{固定体积升温的热容结构。} 任意共存状态的 \(n_v,n_l\) 满足
\(n_v+n_l=N\)。由饱和气体式和克劳修斯---克拉佩龙式， \[
\frac{dn_v}{dT}
=\frac{n_v}{T}\left(\frac{L}{RT}-1\right).
\] 于是固定体积下 \[
\boxed{C_{V,\rm mix}
=n_vC_{V,v}+n_lC_l
+\frac{n_v(L-RT)^2}{RT^2}}.
\]
最后一项是改变相比例的贡献，写成平方后可直接看出非负。这里假设液体体积很小，并用同一温度下的
\(L(T)\) 与热容；有限温区若保留热容差，也应相应更新 \(L(T)\)。

\textbf{接近汽化终点时。} 随温度变化，\(n_v(T)\) 达到 \(N\)
后，潜热项不再存在，后续热容变为
\(NC_{V,v}\)。总内能在相变终点连续，沿过程的导数却可改变。给一个小温升只求微分热量，与一次跨越终点的大温升求积分热量，是两种不同任务。

\needspace{14\baselineskip}

\textbf{给定小热量反求温升。} 原状态的一阶固定体积热容约为
\(C=1330.07\,\mathrm{J/K}\)。若输入
\(100\,\mathrm J\)，且足够小以保持线性化， \[
\Delta T\simeq Q/C\simeq0.075184\,\mathrm K.
\] 相比例变化为 \[
\Delta n_v\simeq0.0326082\,\Delta T
\simeq2.4516\times10^{-3}\,\mathrm{mol}.
\] 饱和压力的相对变化则为 \[
\frac{\Delta p}{p}\simeq\frac{L}{RT^2}\Delta T.
\] 它与 \(\Delta n_v/n_v\) 不相等，因为气体温度也在变： \[
\frac{\Delta n_v}{n_v}
=\frac{\Delta p}{p}-\frac{\Delta T}{T}.
\]
这种反求例子可以检验热量、相比例和压力变化是否在同一阶近似中一致。一次较大加热若改变了相的存在情况，应积分并在相变终点切换模型。

\newpage

\section{第 6 题}\label{ux7b2c-6-ux9898}

在三维空间中给定一个含 \(N\) 个格点（\(N\gg1\)）的自旋为 \(1\)
的体系，每个格点有 \(q\) 个最近邻。格点自旋的 \(z\) 分量 \(S_z\) 可取
\(1,0,-1\)，相应概率分别为 \(P_+,P_0,P_-\)。沿 \(z\) 轴加外磁场 \(h\)
后，系统能量写为 \[
\langle H\rangle=-J\sum_{\langle i,j\rangle}\langle S_iS_j\rangle-h\sum_i\langle S_i\rangle.
\] Bragg--Williams 近似认为不同格点的自旋统计独立，即
\(\langle S_iS_j\rangle=\langle S_i\rangle\langle S_j\rangle\)。设体系均匀磁化，磁化强度定义为
\(m=\langle S_z\rangle=P_+-P_-\)。低温时体系具有非零磁化；温度升高后发生铁磁---顺磁相变，在高于相变温度且
\(h=0\) 时仅有 \(m=0\) 解。

（1）用 \(N,q,J,m,h\) 表示系统的平均内能。

（2）把三种格点数视为概率期望值，利用 \(S=k\ln\Omega\) 求系统熵，其中
\(\Omega\) 为微观状态数。可用 \(\ln(N!)\simeq N(\ln N-1)\)。

（3）在 \(h=0\) 时，利用自由能 \(F=\langle H\rangle-TS\)
取极小值的条件，求 \(P_+,P_0,P_-\) 应满足的方程。

（4）求相变温度 \(T_c\)。

\begin{figure}[H]
\centering
\includegraphics[width=0.719\linewidth,height=\textheight,keepaspectratio,alt={题6　三态自旋格点体系}]{../figures/question-06.pdf}
\caption{题6　三态自旋格点体系}
\end{figure}

\subsection{一、题型定位}\label{ux4e00ux9898ux578bux5b9aux4f4d-5}

本题是统计物理中的平均场近似题。每个格点的自旋变量 \(S_z\) 可取
\(1,0,-1\)，先由统计独立近似写出平均能量，再用多项式计数得到熵，最后极小化自由能，建立序参量
\(m\) 的自洽方程并求临界温度。

\subsection{二、知识点}\label{ux4e8cux77e5ux8bc6ux70b9-5}

\begin{enumerate}
\def\labelenumi{\arabic{enumi}.}
\tightlist
\item
  最近邻相互作用的平均场分解；
\item
  最近邻键数的重复计数；
\item
  多项式系数与 Stirling 公式；
\item
  自由能极小原理与约束变分；
\item
  相变临界点附近的线性化。
\end{enumerate}

\subsection{三、物理图像}\label{ux4e09ux7269ux7406ux56feux50cf-5}

三维晶格中，每个格点有 \(q\) 个最近邻。Bragg--Williams
近似把相邻自旋的关联替换为平均值乘积，从而把多体问题化成单格点概率
\(P_+,P_0,P_-\) 的问题。

\subsection{四、解题入口}\label{ux56dbux89e3ux9898ux5165ux53e3-5}

设 \[
P_++P_0+P_-=1,
\qquad
m=\langle S_z\rangle=P_+-P_-. \tag{1}
\] 能量只需用 \(m\)
表示；熵由三类格点的排列数计算。极小化自由能时不能把三个概率当成完全独立变量，必须同时处理归一化约束。

\subsection{五、解析}\label{ux4e94ux89e3ux6790-5}

\needspace{3\baselineskip}
\textbf{（1）系统的平均内能}

相邻格点统计独立，因此 \[
\langle S_iS_j\rangle
=\langle S_i\rangle\langle S_j\rangle
=m^2. \tag{2}
\] 每个格点有 \(q\) 个最近邻，直接计算得到 \(Nq\)
次，但每一条相邻键会从两个端点各数一次，所以不同相邻键的总数为 \[
N_{\rm bond}=\frac{Nq}{2}. \tag{3}
\] 外场项中有 \(N\) 个格点。于是 \[
\boxed{
\langle H\rangle
=-\frac12NqJm^2-Nhm
}. \tag{4}
\]

\needspace{3\baselineskip}
\textbf{（2）系统熵}

三种自旋的格点数分别为 \[
N_+=NP_+,
\qquad
N_0=NP_0,
\qquad
N_-=NP_-. \tag{5}
\] 在给定三类格点数时，微观状态数是多项式系数 \[
\Omega=\frac{N!}{N_+!N_0!N_-!}. \tag{6}
\] 由 \(S=k\ln\Omega\)，并使用题给近似 \[
\ln n!\simeq n(\ln n-1), \tag{7}
\] 可得 \[
\begin{aligned}
\frac Sk
&\simeq N(\ln N-1)
-\sum_{a=+,0,-}NP_a[\ln(NP_a)-1]\\
&=N\ln N-N
-N\ln N\sum_aP_a
-N\sum_aP_a\ln P_a
+N\sum_aP_a.
\end{aligned}
\] 利用 \(\sum_aP_a=1\)，含 \(N\ln N\) 和常数的项恰好相消： \[
\boxed{
S=-Nk\bigl(P_+\ln P_++P_0\ln P_0+P_-\ln P_-\bigr)
}. \tag{8}
\]

\begin{figure}[H]
\centering
\includegraphics[width=0.537\linewidth,height=\textheight,keepaspectratio,alt={图6-2　自由能极小化与临界点附近的自洽关系}]{../figures/solution-06.pdf}
\caption*{图6-2　自由能极小化与临界点附近的自洽关系}
\end{figure}

\needspace{3\baselineskip}
\textbf{（3）零外场时的平衡概率}

令 \(h=0\)。自由能为 \[
F=-\frac12NqJ(P_+-P_-)^2
+NkT\sum_{a=+,0,-}P_a\ln P_a. \tag{9}
\] 用约束 \[
P_0=1-P_+-P_- \tag{10}
\] 消去 \(P_0\)，再分别对 \(P_+\)、\(P_-\) 求偏导。

先对 \(P_+\) 求导。由于 \(m=P_+-P_-\)，有
\(\partial m/\partial P_+=1\)，故 \[
\frac1N\frac{\partial F}{\partial P_+}
=-qJm+kT\ln\frac{P_+}{P_0}=0. \tag{11}
\] 同理，\(\partial m/\partial P_-=-1\)，于是 \[
\frac1N\frac{\partial F}{\partial P_-}
=qJm+kT\ln\frac{P_-}{P_0}=0. \tag{12}
\] 整理得 \[
\frac{P_+}{P_0}=\exp\!\left(\frac{qJm}{kT}\right),
\qquad
\frac{P_-}{P_0}=\exp\!\left(-\frac{qJm}{kT}\right). \tag{13}
\] 定义 \[
x=\exp\!\left(\frac{qJm}{kT}\right)>0. \tag{14}
\] 则 \[
P_+=xP_0,
\qquad
P_-=\frac{P_0}{x}. \tag{15}
\] 代入归一化条件： \[
P_0\left(x+1+\frac1x\right)=1,
\] 所以 \[
P_0=\frac{x}{x^2+x+1}. \tag{16}
\] 进而 \[
\boxed{
P_+=\frac{x^2}{x^2+x+1},\quad
P_0=\frac{x}{x^2+x+1},\quad
P_-=\frac{1}{x^2+x+1}
}. \tag{17}
\] 另一方面，定义式（14）给出 \[
m=\frac{kT}{qJ}\ln x. \tag{18}
\] 而由三个概率又有 \[
m=P_+-P_-
=\frac{x^2-1}{x^2+x+1}. \tag{19}
\] 因此平衡概率所满足的自洽方程可写成 \[
\boxed{
\frac{kT}{qJ}\ln x
=\frac{x^2-1}{x^2+x+1}
}. \tag{20}
\]

\needspace{3\baselineskip}
\textbf{（4）相变温度}

在临界温度处，非零磁化解从 \(m=0\) 连续产生。令 \[
y=\frac{qJm}{kT},
\qquad x=e^y. \tag{21}
\] 由式（13）也可直接写出 \[
m=\frac{e^y-e^{-y}}{1+e^y+e^{-y}}
=\frac{2\sinh y}{1+2\cosh y}. \tag{22}
\] 临界点附近 \(m\to0\)，因而 \(y\to0\)。使用 \[
\sinh y=y+O(y^3),
\qquad
\cosh y=1+\frac{y^2}{2}+O(y^4), \tag{23}
\] 式（22）的一阶部分为 \[
m\simeq\frac{2y}{3}
=\frac23\frac{qJ}{kT}m. \tag{24}
\] 要出现 \(m\ne0\) 的无穷小解，其一次项系数必须为 \(1\)： \[
1=\frac{2qJ}{3kT_c}. \tag{25}
\] 故 \[
\boxed{T_c=\frac{2qJ}{3k}}. \tag{26}
\]

\textbf{临界条件来自自洽方程斜率。} 把右端记为 \(F(m)\)，顺磁解始终有
\(m=0\)。在原点 \[
F'(0)=\frac{2qJ}{3kT}.
\] \(F'(0)<1\) 时，右端曲线斜率小于直线 \(m\)，零解稳定；降温到
\(F'(0)=1\) 时非零解连续分岔，得到 \(T_c\)。这比直接约去 \(m\)
更清楚：临界点处 \(m=0\)，不能在未展开前随意把原方程除以 \(m\)。

高温极限 \(T\to\infty\)
时三种自旋状态趋于等概率，\(P_+=P_0=P_-=1/3\)、\(m=0\)；低温铁磁相中
\(\pm1\) 态发生不对称占据。两个极限与自洽方程的概率归一化一致。

外场不为零时，\(m=0\)
不再是严格解，临界点被圆滑化为快速交叉；所以式（26）的 \(T_c\)
是零外场下自发对称破缺的温度，不应与有限外场磁化曲线的拐点混为一谈。

\textbf{单点分布与平均场有效哈密顿量。} 用
\(m=\langle S_z\rangle\)，邻点平均贡献为 \(qJm\)。单点有效场
\(H=qJm+h\)，配分函数 \[
Z_1=1+2\cosh(H/kT).
\] 所以 \[
P_0=1/Z_1,\quad P_\pm=e^{\pm H/kT}/Z_1,
\] \[
m=P_+-P_-=\frac{2\sinh(H/kT)}{1+2\cosh(H/kT)}.
\] 其中零自旋状态的统计权重一不可省略；省略后会变成二态 Ising 模型的
\(\tanh\) 式。

\textbf{自由能的双计数修正。}
每个格点用有效场计算时相互作用能被计了两遍，故每点自由能 \[
f(m)=\frac12qJm^2-kT\ln[1+2\cosh((qJm+h)/(kT))].
\] 对 \(m\) 求导： \[
f'(m)=qJ\left[m-\frac{2\sinh y}{1+2\cosh y}\right].
\] 驻点正好给自洽方程，二阶导数决定哪一支稳定。

\textbf{临界展开。} 在 \(h=0\)， \[
\ln(1+2\cosh y)=\ln3+\frac{y^2}3-\frac{y^4}{36}+O(y^6).
\] 所以 \[
f=-kT\ln3+
\left[\frac{qJ}2-\frac{(qJ)^2}{3kT}\right]m^2
+\frac{(qJ)^4}{36(kT)^3}m^4+O(m^6).
\] 二次项在 \(T_c=2qJ/(3k)\)
变号，四次项为正，平均场中发生连续相变。临界下方 \[
\boxed{m^2\simeq\frac83(1-T/T_c)}.
\] 该式只在 \(1-T/T_c\ll1\) 有效，不能令 \(T=0\) 得到超过一的磁化强度。

\textbf{顺磁区线性响应。} 小外场下 \(m\simeq(2/3)(qJm+h)/(kT)\)，所以 \[
\boxed{\chi=\left.\frac{\partial m}{\partial h}\right|_{h=0}
=\frac{2}{3k(T-T_c)}},\qquad T>T_c.
\] 磁化率发散与图6-2中自洽曲线在零点斜率趋一对应。\(J\to0\)
回到单点三态体系 \(\chi=2/(3kT)\)。

\textbf{低温极限。} 取正磁化支，\(H>0\)，则
\(P_+\to1,P_0,P_-\to0\)，\(m\to1\)。能量每点趋
\(-qJ/2-h\)。自旋全部为正时每条邻接键贡献 \(-J\)，总键数
\(Nq/2\)，与能量计数一致。

\textbf{由配分函数求熵与内能。} 设每个格点的平均自旋为 \(m\)，有效场参数
\(y=qJm/(kT)\)。三种自旋概率分别是 \[
p_+=\frac{e^y}{Z},\quad p_0=\frac1Z,\quad
p_-=\frac{e^{-y}}Z,\quad Z=1+2\cosh y.
\] 自洽解满足 \(p_+-p_-=m\)。直接用混合熵计算： \[
s=-k\sum_{\sigma=0,\pm1}p_\sigma\ln p_\sigma
=k\ln Z-kym.
\] 平均场内能要把每条相互作用键只计一次： \[
u=-\frac12qJm^2.
\] 因此 \[
f=u-Ts=\frac12qJm^2-kT\ln Z.
\] 高温解 \(m=0\) 时 \(s=k\ln3\)、\(u=0\)。低温选任一有序支
\(m\to\pm1\)，则 \(s\to0\)、\(u\to-qJ/2\)。若内能写成
\(-qJm^2\)，会把每条键算两遍，且不能与配分函数自由能一致。

\textbf{临界热容的跳跃。} 令
\(\delta=1-T/T_c\)，在临界温度附近展开自由能： \[
f=f_0+a(T)m^2+b(T)m^4+O(m^6),
\] \[
a(T)=\frac{qJ}{2}-\frac{(qJ)^2}{3kT},
\qquad b(T)=\frac{(qJ)^4}{36(kT)^3}.
\] \(kT_c=2qJ/3\)，所以 \[
a=-\frac{qJ}{2}\delta+O(\delta^2),
\qquad b(T_c)=\frac{3qJ}{32}.
\] 极小条件给 \[
m^2=-\frac a{2b}=\frac83\delta+O(\delta^2),
\] \[
f_{\min}-f_0=-\frac{a^2}{4b}
=-\frac23qJ\delta^2+O(\delta^3).
\] 由 \(c=-T\,d^2f_{\min}/dT^2\)，低温支相对高温支的临界热容差为 \[
\boxed{\Delta c=\frac{4qJ}{3T_c}=2k}
\]
（每个格点）。这是本题平均场模型的预测，不能把这个有限跳跃当成任意维数真实系统的精确临界行为。

\textbf{临界等温线的非线性响应。} 增加微小外场能量 \(h\)，有效参数变为
\[
y=\frac{qJm+h}{kT_c}.
\] 利用 \(G(y)=2y/3-y^3/9+O(y^5)\)，在 \(T=T_c\) 留下最低非线性项： \[
m=m+\frac h{qJ}-\frac38m^3+\cdots.
\] 因此 \[
\boxed{m\simeq\left(\frac{8h}{3qJ}\right)^{1/3}}.
\] 临界点不能继续使用线性磁化率，因为该线性系数已发散；三次项决定响应。

\textbf{推广到一般离散自旋。} 作为参数变式，令
\(\sigma=-S,-S+1,\ldots,S\)，各能级在零场时等概率。均方自旋为 \[
\langle\sigma^2\rangle_0=\frac{S(S+1)}3.
\] 线性响应给自洽方程 \[
m\simeq\frac{qJ}{kT}\langle\sigma^2\rangle_0\,m,
\] 所以 \[
\boxed{kT_c=\frac{qJ}{3}S(S+1)}.
\] 配分函数的四阶展开为 \[
\ln Z=\ln(2S+1)+\frac{S(S+1)}6y^2
-\frac{S(S+1)(2S^2+2S+1)}{360}y^4+\cdots.
\] 取 \(S=1\) 恰好得到 \(\ln3+y^2/3-y^4/36\)。若换成只有 \(\pm1\)
两态而没有零态的模型，均方自旋为一，临界温度会改变；因此不能只保留同一个字母
\(m\) 就沿用两态的双曲正切函数。

\needspace{14\baselineskip}

\textbf{直接求解自洽概率的数值例。} 选能量单位 \(qJ=1\)、\(k=1\)，则
\(T_c=2/3\)。在 \(T=0.8T_c\)，解 \(m=2\sinh(m/T)/[1+2\cosh(m/T)]\)
的正根，得到 \[
m\simeq0.677658,\qquad
(P_+,P_0,P_-)\simeq(0.735602,0.206454,0.057944).
\] 三概率之和为1，其差恢复同一个
\(m\)。这时零根依然满足自洽方程，但每格点自由能 \[
f(m)=\frac12m^2-T\ln[1+2\cosh(m/T)]
\] 在非零根比零根低约0.025890，零根不是稳定平衡。在
\(T=1.2T_c\)，零场稳定解为 \(m=0\)，三概率均为 \(1/3\)。

数值求根时，不能用含零点的区间直接宣称找到了铁磁支。先从临界展开估计非零根，再在
\(0<m<1\)
中夹住它；最后核对概率、自洽残差和自由能。若仅要求算法残差小，低温下返回零根也可能``通过''，却没有回答稳定态的问题。

\textbf{增加单离子各向异性的变式。} 明确另加每格点能量 \(D S_z^2\)，其中
\(D>0\) 偏好零自旋态。交换项保持原题不变。记
\(w=e^{-D/(kT)}\)、\(y=(qJm+h)/(kT)\)，单点配分函数和自洽方程为 \[
Z=1+2w\cosh y,\qquad
m=\frac{2w\sinh y}{1+2w\cosh y}.
\] 零场连续相变的候选边界由原点斜率等于1给出： \[
kT_c=\frac{2qJ}{e^{D/(kT_c)}+2}.
\] 例如 \(D/(qJ)=0.10\) 时，数值解为
\(kT_c/(qJ)\simeq0.630545\)，低于原题的 \(2/3\)。让
\(D\to0\)，公式回到原题；若只保留正、负两态，却仍使用相同分母，会错误地忽略零态对配分函数的贡献。

继续展开右端： \[
G(y)=\frac{2w}{1+2w}y
+\frac{w(1-4w)}{3(1+2w)^2}y^3+O(y^5).
\] 当线性项达到临界，且三次项同时为零，得到
\(w=1/4\)，所以候选三临界参数为 \[
kT_t=qJ/3,\qquad D_t=(qJ/3)\ln4.
\] 此时
\(G(y)=y/3-y^5/540+O(y^7)\)，五次项才控制最先出现的非线性回复。原题
\(D=0\)
的连续相变推导不能照搬到整个新参数平面；远离连续边界时须比较完整自由能的极小值，不能仅凭原点斜率决定一阶相变位置。

\newpage

\section{第 7 题}\label{ux7b2c-7-ux9898}

一根质量分布未知的细杆初始竖直立在地面上，接触点为 \(O\)，质量为
\(m\)，质心到 \(O\) 点的距离为 \(r_c\)，杆绕 \(O\) 点的转动惯量为
\(\lambda mr_c^2\)。受微扰后杆开始倒下，重力加速度为 \(g\)。

（1）对均匀杆求 \(\lambda\)。

（2）对任意形状的杆求 \(\lambda\) 的最小值 \(\lambda_{\min}\)。

（3）杆转过角度 \(\theta\)，且 \(O\)
点既未离开地面也未滑动，求此时角速度 \(\dot\theta\)。

（4）若 \(O\) 点铰接，求重力沿杆分量恰好等于质心绕 \(O\)
转动所需向心力时的角度 \(\theta_1\)。

（5）若 \(O\) 点铰接，求质心竖直加速度为零时的角度 \(\theta_2\)。

（6）若 \(O\) 点铰接，求质心水平加速度为零时的角度 \(\theta_3\)。

（7）若 \(O\)
点没有铰接，杆只放在水平地面上，地面支持力保持竖直。存在临界值
\(\lambda_c\)，当 \(\lambda>\lambda_c\) 且摩擦系数足够大时，\(O\)
点一定不会在杆倒地前离开地面，求 \(\lambda_c\)。

（8）在第（7）问情形下取
\(\lambda=1\)，求杆在落地前始终不滑动所需的最小摩擦系数 \(\mu_1\)。

（9）在第（7）问情形下取
\(\lambda=2\)，求杆在落地前始终不滑动所需的最小摩擦系数 \(\mu_2\)。

\begin{figure}[H]
\centering
\includegraphics[width=0.8\linewidth,height=\textheight,keepaspectratio,alt={题7　非均匀细杆由竖直位置倒下}]{../figures/question-07.pdf}
\caption{题7　非均匀细杆由竖直位置倒下}
\end{figure}

\subsection{一、题型定位}\label{ux4e00ux9898ux578bux5b9aux4f4d-6}

本题是刚体平面运动、能量守恒、接触约束与静摩擦条件的综合题。杆的质量分布未知，全部转动性质由无量纲参数
\(\lambda\)
概括。前六问讨论固定端铰接；后三问讨论粗糙水平面上的接触、脱离和不滑动条件。

\subsection{二、知识点}\label{ux4e8cux77e5ux8bc6ux70b9-6}

\begin{enumerate}
\def\labelenumi{\arabic{enumi}.}
\tightlist
\item
  转动惯量与质心位置的矩关系；
\item
  定轴转动的机械能守恒；
\item
  质心的切向与向心加速度；
\item
  接触力、支持力及脱离条件；
\item
  静摩擦不等式的全过程最大值；
\item
  Cauchy--Schwarz 不等式给出的惯量下界。
\end{enumerate}

\subsection{三、物理图像}\label{ux4e09ux7269ux7406ux56feux50cf-6}

以初始竖直向上方向为角度零点，杆倒下时与该方向夹角为 \(\theta\)。质心
\(C\) 到接触点 \(O\) 的距离为 \(r_c\)，绕 \(O\) 的转动惯量为 \[
I_O=\lambda mr_c^2.
\]

\subsection{四、解题入口}\label{ux56dbux89e3ux9898ux5165ux53e3-6}

铰接阶段先由能量求 \(\dot\theta^2\)，再由转动方程求
\(\ddot\theta\)。质心坐标 \[
x_C=r_c\sin\theta,
\qquad
y_C=r_c\cos\theta
\]
统一生成水平、竖直和沿杆方向的加速度。自由接触阶段仍在未滑动时绕瞬时静止点
\(O\)
转动，但支持力和摩擦力必须由质心平动方程求出，并在整个允许角区间检查约束。

\subsection{五、解析}\label{ux4e94ux89e3ux6790-6}

\needspace{3\baselineskip}
\textbf{（1）均匀细杆的 \(\lambda\)}

设杆长为 \(l\)。均匀细杆质心在中点，故 \[
r_c=\frac l2. \tag{1}
\] 绕端点的转动惯量为 \[
I_O=\frac13ml^2. \tag{2}
\] 由定义 \(I_O=\lambda mr_c^2\)： \[
\lambda
=\frac{I_O}{mr_c^2}
=\frac{ml^2/3}{m(l/2)^2}
=\frac43. \tag{3}
\] 所以 \[
\boxed{\lambda=\frac43}. \tag{4}
\]

\needspace{3\baselineskip}
\textbf{（2）任意质量分布时 \(\lambda\) 的最小值}

把杆看作分布在距 \(O\) 为 \(s\) 的位置上的许多质量元 \(dm\)： \[
mr_c=\int s\,dm,
\qquad
I_O=\int s^2\,dm. \tag{5}
\] Cauchy--Schwarz 不等式给出 \[
\left(\int s\,dm\right)^2
\leq\left(\int dm\right)\left(\int s^2\,dm\right)
=mI_O. \tag{6}
\] 代入式（5）： \[
m^2r_c^2\leq mI_O,
\] 于是 \[
I_O\geq mr_c^2,
\qquad
\boxed{\lambda_{\min}=1}. \tag{7}
\] 等号要求所有质量元到 \(O\) 的距离都相同，即质量集中在距 \(O\) 为
\(r_c\) 的位置；一般连续细杆只能无限逼近这一极限。

\needspace{3\baselineskip}
\textbf{（3）转过 \(\theta\) 时的角速度}

杆由静止开始倒下。质心下降高度 \[
\Delta h=r_c-r_c\cos\theta
=r_c(1-\cos\theta). \tag{8}
\] 重力势能减少量转化为定轴转动动能： \[
mgr_c(1-\cos\theta)
=\frac12I_O\dot\theta^2. \tag{9}
\] 代入 \(I_O=\lambda mr_c^2\) 并约去 \(mr_c\)： \[
g(1-\cos\theta)
=\frac12\lambda r_c\dot\theta^2,
\] 故 \[
\boxed{
\dot\theta=\sqrt{\frac{2g(1-\cos\theta)}{\lambda r_c}}
}. \tag{10}
\] 取正根表示 \(\theta\) 随倒下过程增大。

\needspace{3\baselineskip}
\textbf{运动学准备}

重力对 \(O\) 的力矩为 \(mgr_c\sin\theta\)，定轴转动方程给出 \[
mgr_c\sin\theta=I_O\ddot\theta,
\] 所以 \[
\boxed{
\ddot\theta=\frac{g\sin\theta}{\lambda r_c}
}. \tag{11}
\] 质心坐标为 \[
x_C=r_c\sin\theta,
\qquad
y_C=r_c\cos\theta. \tag{12}
\] 连续求两次导数： \[
\begin{aligned}
\ddot x_C
&=r_c(\ddot\theta\cos\theta-\dot\theta^2\sin\theta),\\
\ddot y_C
&=-r_c(\ddot\theta\sin\theta+\dot\theta^2\cos\theta).
\end{aligned} \tag{13}
\] 代入式（10）、（11），记 \(c=\cos\theta\)： \[
\boxed{
\ddot x_C=\frac g\lambda\sin\theta(3c-2)
}, \tag{14}
\] \[
\boxed{
\ddot y_C=-\frac g\lambda(1+2c-3c^2)
}. \tag{15}
\]

\begin{figure}[H]
\centering
\includegraphics[width=0.8\linewidth,height=\textheight,keepaspectratio,alt={图7-2　质心切向、向心加速度及地面对杆的作用力}]{../figures/solution-07.pdf}
\caption*{图7-2　质心切向、向心加速度及地面对杆的作用力}
\end{figure}

\needspace{3\baselineskip}
\textbf{（4）重力沿杆分量等于向心力时的角度}

重力沿杆指向 \(O\) 的分量为 \[
mg\cos\theta. \tag{16}
\] 质心向心加速度大小为 \(r_c\dot\theta^2\)，所需向心力为 \[
mr_c\dot\theta^2. \tag{17}
\] 令二者相等并代入式（10）： \[
mg\cos\theta
=m r_c\frac{2g(1-\cos\theta)}{\lambda r_c}. \tag{18}
\] 约去 \(mg\)，记 \(c=\cos\theta\)： \[
c=\frac2\lambda(1-c).
\] 整理得 \[
(\lambda+2)c=2,
\] 所以 \[
\boxed{
\theta_1=\arccos\frac{2}{\lambda+2}
}. \tag{19}
\]

\needspace{3\baselineskip}
\textbf{（5）质心竖直加速度为零时的角度}

由式（15），令 \(\ddot y_C=0\)： \[
1+2c-3c^2=0. \tag{20}
\] 因式分解： \[
(1-c)(3c+1)=0. \tag{21}
\] 代数方程的两根分别为 \[
c=1,
\qquad
c=-\frac13. \tag{22}
\] 杆从竖直位置倒至水平地面之前，角度范围是 \[
0\leq\theta\leq\frac\pi2,
\qquad 0\leq c\leq1. \tag{23}
\] 因此无地面阻挡的形式铰接模型另有非零解
\(\theta_2=\arccos(-1/3)\)。它位于穿过水平位置之后；若按杆落到水平地面的实际范围，唯一有效解是初始位置
\[
\boxed{\theta_2=0}. \tag{24}
\]

\needspace{3\baselineskip}
\textbf{（6）质心水平加速度为零时的角度}

由式（14）： \[
\sin\theta(3\cos\theta-2)=0. \tag{25}
\] 除去初始位置 \(\sin\theta=0\)，得到 \[
3\cos\theta-2=0,
\] 故 \[
\boxed{
\theta_3=\arccos\frac23
}. \tag{26}
\]

\needspace{3\baselineskip}
\textbf{（7）接触点不离地的临界 \(\lambda_c\)}

现在 \(O\)
不是铰点，而是杆与粗糙地面的接触点。只要尚未滑动且未离地，前面的瞬时定轴运动关系仍成立。竖直方向平动方程为
\[
N-mg=m\ddot y_C. \tag{27}
\] 利用式（15）： \[
N
=mg\left[1-\frac{1+2c-3c^2}{\lambda}\right]. \tag{28}
\] 全过程不离地要求对运动经过的所有 \(c\) 都有 \(N\geq0\)，即 \[
\lambda\geq1+2c-3c^2. \tag{29}
\] 右端是开口向下的二次函数。配方： \[
1+2c-3c^2
=\frac43-3\left(c-\frac13\right)^2. \tag{30}
\] 最大值为 \(4/3\)，在 \(c=1/3\) 处取得。因此 \[
\boxed{\lambda_c=\frac43}. \tag{31}
\] 当 \(\lambda>4/3\)
时，只要静摩擦足够，支持力在倒地前始终非负；\(\lambda<4/3\)
时，支持力会先降到零，原接触约束失效。

\textbf{（8）\(\lambda=1\) 的接触阶段与全过程。} 此时 \[
N=mgc(3c-2),\qquad
f=mg\sin\theta(3c-2).
\] 接触可维持的初段为 \(c>2/3\)。在这段， \[
\frac{|f|}{N}=\tan\theta.
\] 所以保持接触直至自然脱离而不先发生滑动，所需摩擦系数为 \[
\boxed{\mu_{\rm before\ detach}\ge\sqrt5/2}.
\] 在 \(c=2/3\) 时支持力和摩擦力同时为零；之后形式固定端解要求
\(N<0\)，地面无法提供向下拉力。若题意要求落地前始终保持 \(O\)
固定且贴地，则 \[
\boxed{\text{不存在可保证全过程固定接触的有限摩擦系数}}.
\] \(\sqrt5/2\)
只保证脱离之前不滑动，不能保证脱离之后继续绕固定端转动。原第（7）问的持续接触前提
\(\lambda>4/3\) 与本问 \(\lambda=1\) 不能同时成立，应明确这一题设边界。

\needspace{3\baselineskip}
\textbf{（9）\(\lambda=2\) 时不滑动所需的最小摩擦系数}

此时式（27）给出 \[
N=\frac{mg}{2}(1-2c+3c^2). \tag{38}
\] 二次式 \(1-2c+3c^2\) 的判别式为 \((-2)^2-12=-8<0\)，所以 \(N\)
始终为正，杆在倒地前不会脱离。

由式（31）、（37）： \[
\frac{|f|}{N}
=\frac{\sqrt{1-c^2}\,|3c-2|}{1-2c+3c^2}
\equiv R(c),
\qquad 0\leq c\leq1. \tag{39}
\] 最小摩擦系数应覆盖全过程，因此 \[
\mu_2=\max_{0\leq c\leq1}R(c). \tag{40}
\] 为避免绝对值和根式，可在两个分段上极大化 \(R^2\)： \[
R^2(c)=\frac{(1-c^2)(3c-2)^2}{(1-2c+3c^2)^2}. \tag{41}
\] 求导并筛选区间内驻点，决定全局最大值的根为 \[
c_*\approx0.07794. \tag{42}
\] 代回式（38）： \[
R(c_*)\approx2.0419.
\] 故三位有效数字为 \[
\boxed{\mu_2=2.04}. \tag{43}
\]

\textbf{极值必须包含区间端点。} \(R(c)\) 含绝对值，先在 \(0\le c\le2/3\)
与 \(2/3\le c\le1\) 分段极大化 \(R^2\)，再比较驻点和端点。端点给 \[
R(0)=2,\qquad R(2/3)=0,\qquad R(1)=0.
\] 内部根 \(c_*\simeq0.07794\) 给 \(2.0419\)，只略大于
\(R(0)\)。若只用粗略作图，很容易误把水平位置对应的 \(R(0)=2\)
当成最大值。

支持力分母 \(1-2c+3c^2\)
始终正，说明整个过程中接触保持；摩擦比最大值有限，所以只要
\(\mu\ge2.04\)，无滑动假设自洽。

变量 \(c=\cos\theta\)
随杆从竖直到水平由1降到0，因此``首先到达''的物理解要按 \(c\)
递减顺序筛选，而不是按代数根从小到大排列。涉及支持力为零或摩擦达限的多根方程时，这个运动方向条件决定哪一个根先发生。

\textbf{通用初倾角。} 正文从竖直静止的极限开始。若初始角为
\(\theta_i\)，静止释放，则固定端模型能量 \[
\frac12\lambda mr_c^2\dot\theta^2
=mgr_c(\cos\theta_i-\cos\theta).
\] 记 \(c_i=\cos\theta_i\)、\(c=\cos\theta\)，有 \[
\dot\theta^2=\frac{2g(c_i-c)}{\lambda r_c},
\qquad
\ddot\theta=\frac{g\sin\theta}{\lambda r_c}.
\] 质心坐标为 \(x_C=r_c\sin\theta,y_C=r_c\cos\theta\)。二次求导： \[
a_x=\frac g\lambda\sin\theta(3c-2c_i),
\] \[
a_y=-\frac g\lambda(1+2c_ic-3c^2).
\] 令 \(c_i=1\)
回到正文。这个变式可检查重力切向加速度和向心加速度的两个分量是否画反。

\textbf{持续接触的临界值。} 支持力为 \[
N=mg\left[1-\frac{1+2c_ic-3c^2}{\lambda}\right].
\] 在全过程 \(0\le c\le c_i\) 内，二次函数顶点为 \(c=c_i/3\)，最大值 \[
1+c_i^2/3.
\] 因此 \[
\boxed{\lambda>1+c_i^2/3}
\] 保证正支持力。等号时在一个瞬间支持力为零，只能称接触边界；竖直极限给
\(4/3\)。同时还要静摩擦足够，支持力非负本身不保证端点不滑动。

\textbf{脱离角。} 若 \(\lambda\) 小于临界但初态接触可维持，解 \(N=0\)：
\[
3c^2-2c_ic+\lambda-1=0.
\] 沿 \(c\) 逐渐减小的路径应取首次遇到的较大根： \[
c_{\rm detach}=\frac{c_i+\sqrt{c_i^2-3(\lambda-1)}}3.
\] 根必须在
\([0,c_i]\)；若在外面，说明该参数分支或初态假设不满足。竖直、\(\lambda=1\)
得 \(c=2/3\)。

\textbf{为什么竖直静止释放仍会倒。} 精确理想竖直态
\(\theta=0,\dot\theta=0\)
是不稳定平衡，若没有扰动，数学上会永远停在那里。题目写``由于微扰倒下''，实际取
\(\theta_i\downarrow0\) 的极限。小角阶段 \[
\ddot\theta\simeq\frac g{\lambda r_c}\theta.
\] 令 \(\omega_*=\sqrt{g/(\lambda r_c)}\)： \[
\theta(t)=\theta_i\cosh(\omega_*t).
\] 到达固定有限角的等待时间随 \(\theta_i\to0\)
对数增大；角度之间的能量关系有有限极限，绝对倒下时间却没有无扰动的有限值。

\textbf{从竖直位置释放的时间问题。}
正文的角速度与力适用于杆已经开始倒下后的运动。完全竖直且角速度严格为零的理想杆处于不稳定平衡，运动方程仍允许它永远保持竖直。因此``从竖直静止开始倒下''实际指一个很小的初倾角极限。

取原能量式对应的竖直极限轨道，记 \[
D=L^2+L_0^2,\qquad
\dot\theta^2=\frac{2gL}{D}(1-\cos\theta).
\] 从小角 \(\varepsilon>0\) 到 \(\theta\) 的时间为 \[
t(\theta)-t(\varepsilon)
=\sqrt{\frac{D}{2gL}}\int_\varepsilon^\theta
\frac{d\vartheta}{\sqrt{1-\cos\vartheta}}.
\] 由于 \(1-\cos\vartheta=2\sin^2(\vartheta/2)\)， \[
\int\frac{d\vartheta}{\sqrt{1-\cos\vartheta}}
=\sqrt2\ln\tan\frac{\vartheta}{4}.
\] 所以 \[
\boxed{t(\theta)-t(\varepsilon)
=\sqrt{\frac D{gL}}
\ln\frac{\tan(\theta/4)}{\tan(\varepsilon/4)}}.
\] \(\varepsilon\to0\)
时该时间对数发散，而到达固定角度后的角速度有有限极限。能量解能描述相空间轨迹，但不能凭空确定从不稳定平衡点离开的时刻。

\textbf{小角指数增长。} 对小 \(\theta\)，动力学方程给 \[
\ddot\theta=\frac{gL}{D}\theta+O(\theta^3).
\] 令 \(\nu=\sqrt{gL/D}\)。若真实初始条件是
\(\theta(0)=\varepsilon\)、\(\dot\theta(0)=0\)，则 \[
\theta(t)\simeq\varepsilon\cosh(\nu t).
\] 达到仍属小角范围的 \(\theta_f\gg\varepsilon\) 时， \[
t_f\simeq\frac1\nu\operatorname{arcosh}\frac{\theta_f}{\varepsilon}
\simeq\frac1\nu\ln\frac{2\theta_f}{\varepsilon}.
\] 这与竖直极限轨道的对数规律一致，但两者的初始能量相差
\(O(\varepsilon^2)\)，不应把两种时间公式混成完全相同的初值问题。

\textbf{接触力必须随时间检查。} 知道 \(\theta(t)\) 后，地面支持力为 \[
N(t)=mg\left[
1-\frac{1+2\cos\theta(t)-3\cos^2\theta(t)}{\lambda}
\right].
\] 在 \(N>0\) 的阶段才能继续使用固定端点运动。若出现
\(N=0\)，下一阶段应改用自由抛转；若先达到
\(|f|=\mu N\)，下一阶段应改用滑动接触。时间参数化并不会延长固定端点模型的有效区间。

\needspace{14\baselineskip}

\textbf{法向力和摩擦力同时为零的特殊点。} 对
\(\lambda=1\)，由固定端点运动方程 \[
N=mg\,c(3c-2),\qquad
f=mg\sqrt{1-c^2}(3c-2).
\] 在仍接触的区间 \(2/3<c\le1\)， \[
\frac{|f|}{N}=\frac{\sqrt{1-c^2}}c.
\] 于是当 \(c\downarrow2/3\)， \[
\mu_{\rm req}\to\frac{\sqrt5}{2}.
\] 这里两个反力含相同的一阶因子，约去后比值有有限极限。与
\(\lambda=4/3,c=1/3\) 的情形不同，后者 \(N\) 二阶为零而 \(f\)
不为零，比值发散。接触失效点不能只看分子或分母之一，必须比较完整表达式和到达该点之前的支路。

\newpage

\section{第 8 题}\label{ux7b2c-8-ux9898}

质量均为 \(m\)
的绝缘均匀圆环与光滑带电珠，在足够粗糙的水平地面上运动。环半径
\(r\)，珠电量 \(q\)，磁场 \(B\) 垂直纸面向里。\(\beta\)
为圆心到珠的半径与向下竖直方向的有向夹角。初态按原图取珠在左侧点
\(\beta_0=\pi/2\)，两者静止，环始终纯滚动并不离地。定义
\(\lambda=q^2B^2r/(m^2g)\)。

（1）求再次相对静止的非零位置开始出现的
\(\lambda_c\)。（2）\(\lambda>\lambda_c\) 时求 \(\cos\beta_c\)。
（3）（4）珠经过 \(\beta_1\)
时求相对圆环角速度的较小、较大大小。（5）（6）求相应圆心速度的两个大小。后四问取
\(qB/m=1\,\mathrm{s^{-1}},r=1\,\mathrm m,g=10\,\mathrm{m\,s^{-2}}\)；原题未给
\(\beta_1\) 数值，答案保留该参数。

\includegraphics[width=0.86\linewidth,height=\textheight,keepaspectratio,alt={题8　题设与符号}]{../figures/question-08.pdf}

\subsection{一、题型定位}\label{ux4e00ux9898ux578bux5b9aux4f4d-7}

纯滚动圆环与带电珠的耦合运动。角度相对空间定义，题目后两问却要求相对环的角速度，二者相差环的自转速度。

\subsection{二、知识点}\label{ux4e8cux77e5ux8bc6ux70b9-7}

圆环纯滚动；珠的空间速度与相对速度；磁场不做功；磁冲量与水平动量；耦合动能的正定性。

\subsection{三、物理图像}\label{ux4e09ux7269ux7406ux56feux50cf-7}

图8-1中磁场向里，\(x\) 向右；\(\beta\)
从竖直向下方向转向左侧。初态珠在左侧点，即
\(\beta_0=\pi/2\)。这个初态位置须由原图补充。圆环顺时针角速度为
\(u=\dot x/r\)；珠的空间角速度为 \(w=\dot\beta\)，相对环角速度为
\(w-u\)。

\subsection{四、解题入口}\label{ux56dbux89e3ux9898ux5165ux53e3-7}

先求动能和磁冲量，再消去 \(u\)。原题只给 \(\beta_1\)，没有给
\(\beta_1=1\)；因此保留 \(c_1=\cos\beta_1\)。取一弧度只能作为附加算例。

\subsection{五、解析}\label{ux4e94ux89e3ux6790-7}

\textbf{运动学与能量。} 珠的位置和速度为 \[
x_P=x-r\sin\beta,\qquad y_P=r-r\cos\beta,
\] \[
\dot x_P=\dot x-rw\cos\beta,\qquad \dot y_P=rw\sin\beta.
\] 圆环平动与转动的总动能是 \(m\dot x^2\)。珠动能展开后，系统能量守恒为
\[
\frac32u^2-uwc+\frac12w^2=\frac grc,\qquad c=\cos\beta.
\] 这里右端是珠从侧点下降所释放的重力势能，已除以 \(mr^2\)。

\textbf{磁冲量与滚动。} 珠受 \(F_{Bx}=-qB\dot y_P\)；其水平冲量为
\(qBr c\)。地面摩擦冲量使环纯滚动，等于 \(-m\dot x\)。因此 \[
3u-cw=\omega_Bc,\qquad \omega_B=qB/m.
\] 这是一阶积分，不是水平机械动量守恒。代入能量，交叉的一次项相消： \[
(3-c^2)w^2+\omega_B^2c^2=6gc/r.
\] 所以 \[
w_\pm=\pm W(c),\qquad
W(c)=\sqrt{\frac{6gc/r-\omega_B^2c^2}{3-c^2}},
\quad
u_\pm=\frac{c(\pm W+\omega_B)}3.
\]
初始下落对应负号支，回升对应正号支；给定角位置可能先后经过两次。根号非负与
\(0\le c\le1\) 共同限制可达位置。

\textbf{再次相对静止。} 珠随环材料点一起转动时 \(w=u\)，而不是
\(w=0\)。对非零 \(c\) 联立得 \[
(3-c)u=\omega_Bc,\qquad (2-c)u^2=gc/r.
\] 因此 \[
\lambda=\frac{\omega_B^2r}{g}
=\frac{(3-c)^2}{(2-c)c}.
\] 右式在 \(0<c\le1\) 随 \(c\) 减小而增大，最小值为 \(4\)。故 \[
\boxed{\lambda_c=4}.
\] 非零物理解满足 \[
(\lambda+1)c^2-(2\lambda+6)c+9=0,
\] \[
\boxed{c=\frac{\lambda+3-\sqrt{\lambda^2-3\lambda}}{\lambda+1}},\qquad \lambda>4.
\] 判别式为 \(4\lambda(\lambda-3)\)。\(\lambda\ge3\)
仅保证代数实根；\(3\le\lambda<4\) 的根仍在物理区间之外。零根 \(c=0\)
不能在约去 \(c\) 后遗失。

\textbf{原题给定 \(\beta_1\) 的答案。} 代入
\(\omega_B=1\,\mathrm{s^{-1}}\)、\(r=1\,\mathrm m\)、\(g=10\,\mathrm{m\,s^{-2}}\)：
\[
W_1=\sqrt{\frac{60c_1-c_1^2}{3-c_1^2}}\,\mathrm{s^{-1}}.
\] 在可达区间且 \(c_1>0\)，相对角速度的两个大小为 \[
\boxed{|\omega_{\rm rel}|_{\rm small}
=\frac{|(3-c_1)W_1-c_1|}{3}},
\] \[
\boxed{|\omega_{\rm rel}|_{\rm large}
=\frac{(3-c_1)W_1+c_1}{3}}.
\] 圆心速度大小为 \[
\boxed{|\dot x|_{\rm small}=\frac{c_1|W_1-1|}3\,\mathrm m},
\qquad
\boxed{|\dot x|_{\rm large}=\frac{c_1(W_1+1)}3\,\mathrm m}.
\] 上式中的 \(W_1\) 数值按秒的单位代入；更一般时应写
\(r c_1|W_1-\omega_B|/3\) 和 \(r c_1|W_1+\omega_B|/3\)。本题
\(\lambda=0.1<4\)，\(W_1>1\) 的适用区间是
\(60c_1>3\)，若更靠近侧点则须保留绝对值重新排序，不能把这组大小标签延伸到整个区间。

\textbf{附加算例 \(\beta_1=1\,\mathrm{rad}\)。} 此时
\(c_1=0.540302306\)， \[
W_1=3.444288,\qquad
u_-=-0.440218,\quad u_+=0.800420.
\] 于是相对角速度大小为 \[
\boxed{2.643868,\ 3.004070\ \mathrm{s^{-1}}},
\] 圆心速度大小为
\(0.440218,\ 0.800420\,\mathrm{m\,s^{-1}}\)。两支空间角速度大小相同，相对角速度大小不同；不能把
\(3.444288\) 填入相对速度答案。

\begin{figure}[H]
\centering
\includegraphics[width=0.86\linewidth,height=\textheight,keepaspectratio,alt={图8-2　推导所用的几何与方向}]{../figures/solution-08.pdf}
\caption*{图8-2　推导所用的几何与方向}
\end{figure}

\textbf{一般环与珠质量比。} 珠质量为 \(m\)，环质量为 \(M\)，环的惯量为
\(I=\kappa Mr^2\)。定义 \[
A=1+(1+\kappa)M/m.
\] 初态仍为侧点静止，能量与冲量积分分别变为 \[
\frac A2u^2-uwc+\frac12w^2=gc/r,
\qquad Au-cw=\omega_Bc.
\] 均匀环且 \(M=m\) 给 \(A=3\)。消去 \(u\)： \[
(A-c^2)w^2+\omega_B^2c^2=2Agc/r.
\] 所以 \[
w_\pm=\pm\sqrt{\frac{2Agc/r-\omega_B^2c^2}{A-c^2}},
\quad u_\pm=c(w_\pm+\omega_B)/A.
\] 动能正定要求 \(A-c^2>0\)，对正质量与正惯量自动满足。

\textbf{通用相对静止条件。} \(w=u\) 给 \[
(A-c)u=\omega_Bc,\qquad
\frac{A+1-2c}{2}u^2=gc/r.
\] 因此 \[
\lambda=\frac{2(A-c)^2}{c(A+1-2c)}.
\] 在本题 \(A=3\) 回到原式。一般临界是否总在 \(c=1\)
需对该函数求导，不可把特殊数值四直接移植。对 \(A>1\) 导数分子正比于
\((A-c)[(3A-1)c-A(A+1)]\)，还要检查其零点是否在允许区间。

\textbf{空间转向位置。} 空间半径转向满足 \(w=0\)。除 \(c=0\)
外，能量根号给 \[
c_{\rm turn}=2A/\lambda.
\] 当 \(\lambda<2A\)，这根在一之外，空间转向要在另一个侧点；当
\(\lambda>2A\)，珠在到达底部前就出现空间转向。相对静止出现的临界与空间转向临界不相同：本题分别为四与六。

\textbf{绝对值排序的完整形式。} 对给定可达 \(c\)，令
\(W>0\)。两支相对速度是 \[
\omega_{\rm rel,\pm}=
\pm(1-c/A)W-\frac{\omega_Bc}{A}.
\] 它们的大小必须写成 \[
\left|(1-c/A)W-\omega_Bc/A\right|,
\quad
\left|(1-c/A)W+\omega_Bc/A\right|.
\] 圆心速度大小为 \(r c|W-\omega_B|/A\) 与
\(r c|W+\omega_B|/A\)。取正磁场且 \(c>0\)
时后者较大，但前者的内部差不一定为正。原题参数下侧点附近
\(W<1\)，不能省绝对值。

\textbf{用动能配方检验磁场能量记账。} 在一般质量与转动惯量参数 \(A\)
下，单位 \(mr^2\) 的动能是 \[
\mathcal T=\frac A2u^2-cuw+\frac12w^2,\qquad c=\cos\beta.
\] 配方得 \[
\mathcal T=\frac A2\left(u-\frac cA w\right)^2
+\frac12\left(1-\frac{c^2}{A}\right)w^2.
\] 因为
\(A>1\)、\(|c|\le1\)，动能正定。磁场虽然改变平动与转动的分配，但不向系统输入机械能。由冲量关系
\[
Au-cw=\omega_Bc
\] 可把配方中的第一项化成 \(\omega_B^2c^2/(2A)\)，能量式给 \[
w^2=\frac{2Agc/r-\omega_B^2c^2}{A-c^2}.
\] 分母恒正，因此能否到达某个角度只取决于分子非负： \[
2Ag/r-\omega_B^2c\ge0\quad(c>0).
\] 这是粒子在空间中的转向条件，和相对圆环静止的条件 \(w=u\) 不同。

\textbf{圆环位移的积分表达。} 沿一条不经过 \(w=0\) 的运动支路， \[
\frac{dx}{d\beta}=\frac{\dot x}{\dot\beta}
=\frac{rc}{A}\left(1+\frac{\omega_B}{w}\right).
\] 从初始 \(\beta_0=\pi/2\) 出发： \[
x(\beta)-x_0=
\frac rA(\sin\beta-1)
+\frac{r\omega_B}{A}
\int_{\pi/2}^{\beta}\frac{\cos b}{w(b)}\,db.
\] 无磁场时第二项消失，得到精确关系 \[
\boxed{x(\beta)-x_0=\frac rA(\sin\beta-1)}.
\] 到最低点 \(\beta=0\) 时圆环向左移动
\(r/A\)。这可直接用水平动量守恒检验：粒子相对环心的水平坐标由 \(-r\)
增到零，圆环必须反向移动以保持系统水平质心不动。

\textbf{初始点的时间积分是否发散。} 在 \(\beta\) 接近 \(\pi/2\) 时，记
\(\epsilon=\pi/2-\beta\)，则 \(c\simeq\epsilon\)， \[
|w|\simeq\sqrt{\frac{2g}{r}\epsilon}.
\] 因此 \[
\int_{\beta}^{\pi/2}\frac{db}{|w(b)|}
\simeq\sqrt{\frac{2r}{g}}\sqrt{\epsilon},
\]
是有限的。与上一题竖直杆的指数离开不同，此处初始切向重力不为零，粒子会立即加速。

若场很强，在内部角度 \(c_t=2Ag/(r\omega_B^2)<1\) 出现
\(w=0\)，速度积分应在转向点分段，改换 \(w\)
的正负支。直接在全程使用同一个平方根符号，会给出错误的时间与圆环位移。

\needspace{14\baselineskip}

\textbf{最低点的空间速度与相对速度。} 原题 \(A=3\)，记
\(\lambda=\omega_B^2r/g\)。在 \(c=1\)，空间角速度满足 \[
w^2=\frac gr\left(3-\frac{\lambda}{2}\right).
\] 最低点可达需 \(\lambda\le6\)。两支分别有 \[
w_\pm=\pm\sqrt{\frac gr(3-\lambda/2)},\qquad
u_\pm=(w_\pm+\omega_B)/3.
\] 相对角速度为 \[
\omega_{{\rm rel},\pm}=
\frac{2w_\pm-\omega_B}{3}.
\] 若 \(\omega_B>0\)，相对静止在最低点要求正支 \(2w_+=\omega_B\)，得到
\[
\boxed{\lambda=4}.
\] 空间转向阈值却为 \(\lambda=6\)。在 \(\lambda=6\) 时 \(w=0\)，但
\(u=\omega_B/3\)，粒子仍相对圆环转动。这个代入检验用最简单的一个几何位置区分了两种``静止''。
\end{document}
